% =====================================================================
%  R. Nielsen, ON THE CONDITIONS OF SOCIETY.  Six Popular Lectures.  Copenhagen: Gyldendal (F. Hegel & Søn), 1877.
%  103 pp.  (Om Samfundets Vilkaar.  Sex folkelige Foredrag.)
%  ENGLISH TRANSLATION, from transcription.tex (Danish).
%  GENERATED by ebook-method/om-samfundets-vilkaar/tr_build.py --write from this template and the six translators'
%  parts in ebook-method/om-samfundets-vilkaar/tr/ (TRANSLATOR-BRIEF.md; terms: GLOSSARY.md, the title and three key
%  terms approved by Hans 2026-10-12).  Conventions: TRANSLATION-PLAYBOOK.md.  Page numbers of the 1877 print in the
%  margin.
% =====================================================================
\documentclass[12pt,a4paper,openany]{book}
\usepackage[utf8]{inputenc}
\usepackage[T1]{fontenc}
\usepackage{libertinus}
\usepackage{libertinust1math}
\usepackage[english]{babel}
\usepackage[protrusion=true,expansion=true]{microtype}
\usepackage{setspace}
\usepackage[margin=1.3in]{geometry}
\usepackage{fancyhdr}
\usepackage{hyperref}
\hypersetup{pdftitle={On the Conditions of Society}, pdfauthor={Rasmus Nielsen}, hidelinks}
\setlength{\headheight}{15pt}
\pagestyle{fancy}
\fancyhf{}
\fancyhead[LE,RO]{\thepage}
\fancyhead[C]{\textit{R. Nielsen: On the Conditions of Society (1877)}}
\setstretch{1.3}
\usepackage{marginnote}
\renewcommand{\thefootnote}{\ifcase\value{footnote}\or*)\or**)\or***)\fi}   % as the print: *) and **), restarting each page
\newcommand{\tocline}[3]{\noindent\makebox[2.2em][r]{#1}\hspace{0.6em}#2\dotfill\ #3\par}
\newcommand{\opage}[1]{\marginnote{\footnotesize\textit{[#1]}}}
% a lecture opens a new page: its numeral and title, centred, then a short rule
\newcommand{\lecture}[2]{\clearpage\phantomsection\addcontentsline{toc}{chapter}{#1 #2}\begin{center}{\large #1}\\[0.6em]{\large #2}\\[0.4em]\rule{3em}{0.4pt}\end{center}\medskip\noindent\ignorespaces}

\begin{document}

% PDF 7, the title page
\begin{titlepage}
\begin{center}
\vspace*{2em}
{\Huge On the Conditions of Society.}\\[1.5em]
{\large Six Popular Lectures}\\[1em]
{\small by}\\[1em]
{\large R. Nielsen.}\\[0.3em]
\rule{2em}{0.4pt}\\[2em]
[publisher's device]\\[4em]
{\large Copenhagen.}\\[0.5em]
Published by the Gyldendal Bookshop (F. Hegel \& Søn).\\[0.5em]
{\small Printed by J. Jørgensen \& Co.}\\[0.5em]
{\small 1877.}\\[3em]
{\small Translated from the Danish}
\end{center}
\end{titlepage}

% PDF 9, the print's Indhold
\begin{center}{\large\emph{Contents}.}\\[0.4em]\rule{3em}{0.4pt}\end{center}
\hfill{\small Page}\par
\tocline{I.}{The Relation of Our Constitution to the Political Movements of the Moment (given in Kolding)}{1.}
\tocline{II.}{Popular Self-Government: People and King (given in Kolding)}{18.}
\tocline{III.}{Political Experiences: Parliamentarism (given in Kolding)}{36.}
\tocline{IV.}{Free Constitution and Freedom (given in Middelfart)}{54.}
\tocline{V.}{Aristocratism and Democratism (given in Viborg)}{72.}
\tocline{VI.}{Politics and Conscience: Machiavellianism (given in Fredericia)}{89.}
\clearpage

% --- p. 1 ---
\setcounter{footnote}{0}%
\lecture{I.} {The Relation of Our Constitution to the Political Movements of the Moment.}
\opage{1}An old fable tells that a general war had once broken out
between the four-footed animals and the birds.
The animal world had divided into two fighting
armies, which advanced: it was a matter of victory or
defeat. Every living creature had taken its side;
yet there was one creature that would not declare itself,
but wished to hold itself hovering above both parties, and that was
the bat. It thought as follows: in troubled times
one must be cautious, not hasty, but look things
over and keep on good terms with everyone. No one can know
to which side victory will incline; the wise man is prepared
for anything: if the four-footed win, then I belong to
their number, for I too am four-footed; but should it
happen that the birds carried off the victory, well then, I
have wings after all and can very well pass for a bird.
Besides, there is something lofty and superior in raising oneself
thus above the parties; one is so unbiased, so
many-sided, so honorably disinterested. And when the passions
subside, what then? Then I become a peacemaker,
recommend mutual accommodation, initiate
agreements and am always on top. Thus the bat
thought; for the bat was a
% --- p. 2 ---
\setcounter{footnote}{0}%
\opage{2}politician, a diplomatically refined politician. But the contending
parties would not be fooled; they soon noticed
the ambiguous creature, and however much at odds they were among themselves,
they nevertheless at once agreed to strike
the ambiguous one dead.

As it goes with the bat in the fable, so, the very timid
think, it must surely go with anyone
who in an agitated time wishes to step forward before the people
with a word of guidance, if, mark you, he does not
himself wish to take sides, but confines himself to attempts
to interpret the movements and to cast stray beams of light upon the
arena of the parties. --- Yet the prospects are not so dark.
To be sure, history knows party struggles in which the passions
cry so loudly that reasonable arguments must
be silent; but these are the terrors of the thunderstorm, earthquakes that overturn society.
Such terrors of the thunderstorm have so far, fortunately,
kept outside the borders of our fatherland. We
are waking from an old drowse, but we shall not on that account
fly up in passion and abandon ourselves to
the blind impulses of unclear moods. We shall think
before we act, and test before we judge. The testimony still
stands firm, the honorable testimony of history, that
the Danish people is a calm, sober, circumspect on all sides
people.

We listen to the words of those who wish to establish peace; but
we do not care for the political bats: they shun the
clear light and prefer to swarm in the twilight. We heed
the voice of prudence, but we cannot admire
the ambiguous go-betweens whose skill it is
to break the point off every principle and to obscure what
lies at the bottom. We value reconciliation and
moderation, but are on our guard against obliging
% --- p. 3 ---
\setcounter{footnote}{0}%
\opage{3}conciliation commissioners who turn up with a stock of all sorts of
opinions, adapted to the circumstances,
and change color like the four-legged creature called the chameleon.
We have an idea that, in spite of their apparent
impartiality, they are deeply entangled in the web of
party interests, that, to put it plainly, they have something
behind the ear. A people never puts trust in those who
give the impression of having something behind the ear. There is a
popular virtue called honesty; to honesty corresponds
openness, and to openness the frankness which, in
seemly self-restraint, speaks out and
gives us the matter as it is, whether it is
otherwise pleasant or unpleasant to us.

By human conditions it is surely impossible that
a living people can move freely under a free constitution
without parties forming. Capable party leaders
are indispensable figures in society; every
party must then stand and fall with the cause for which
it fights. But the popular lecture cannot
take sides; that lies in the nature of the thing, it is impossible.
With whom should it take sides? With the people! ---
But the people as a people is no party. The popular
essence stirs, more or less obscured, in all
parties --- on the Right as well as on the Left. The
popular essence encompasses the whole, ramifies into
everything and pervades us all. The popular element
enters into the most various mixtures; the task
is to get it separated from these mixtures as purely
as possible. When two or more parties are quarreling about a
matter important to the whole people, and the popular yardstick
is applied, it is a matter of course that one party can
be placed in a more favorable light than another; but
% --- p. 4 ---
\setcounter{footnote}{0}%
\opage{4}in this there is no partiality. What matters
is only this: that the yardstick is right, and the party
rightly judged. To present a party either worse or
better than it demonstrably is, is partisan; but to wish, without any
side considerations, to characterize it as it is, is
not partisan.

It will cost labor before the great general public
reaches the point where, unbiased by party considerations, it can
distinguish sharply between the popular and the unpopular.
A gladdening sign has nevertheless already been given. The ordinary
people have begun to open their eyes to a very
important distinction, the distinction between practical politics
and an elucidation of political conditions comprehensible to all. A
small example. No one shall be able with justice to say that
the population of this region lacks a practical sense for
the burning questions of the moment. And no one shall
with justice be able to say that in the lectures which I last
year had the honor of giving from this place before
a numerous circle of benevolent listeners, sides were taken
either for the Right or for the Left. When a wish has
nevertheless now been expressed that the general points of view
which were then put forward might, by a further
continuation, find closer clarification and more distinct
application to given cases, I may surely
without presumption regard the honorably repeated invitation
as a proof that the distinction mentioned
is beginning to make its way through.

The popular lecture can go along very different
paths, pluck fruits from different trees and
draw from many sources. It can fetch its subjects from
mythology, legendary history, general world history % print: Verdenhistorie (misprint)
etc.; but with all this there is one thing that must never
% --- p. 5 ---
\setcounter{footnote}{0}%
\opage{5}be lost from sight. The people lives in the present.
The present, real life has a multiplicity
of conditions, troubles, entanglements and questions
which daily announce themselves with great insistence; to
the answering of these questions the popular lecture must
first and last point, if only indirectly.
The popular lecture is not to strike down upon the momentary
with the assertion that this or that is
the only thing that ought to be done; it is much rather to aim
at elucidating the existing conditions in such a way
that each individual among the people and thereby the whole people gradually
can raise itself to that level of survey and
insight at which the majority and, if possible, all are put in
a position to pass a reasonable, independent and well-grounded
judgment.

How this, as I conceive it, is to be
understood will be seen as I now proceed to
treat the announced subject of the present lecture: the
\emph{relation of our Constitution to the political
movements of the moment}.

The little book with the title: \emph{The Revised
Constitution of the Kingdom of Denmark of 5 June 1849},
is in extent perhaps nearly the smallest, but in
content the greatest in our literature. It is a book of the people
so comprehensible to all that it can be read and understood
without juridical learning. In this there lies, then, a summons
to everyone who wishes to form an independent judgment
about our constitution, not to content himself with extracted
sentences, but to know the law in its coherence
and to make himself familiar with its spirit and letter.

That the Constitution divides the entire power of the state into three
powers, the legislative, the executive and the
% --- p. 6 ---
\setcounter{footnote}{0}%
\opage{6}judicial, in such a way that (§ 2) the legislative power rests with
the King and the Rigsdag jointly, gives us in this
connection no occasion for special consideration.
But the division of the Rigsdag into two awakens our attention.
``The Rigsdag,'' it says (§ 29), ``consists of the Folkething
and the Landsthing.'' What does this division into two mean?
Why could the Folkething alone not have been
enough? The only information we can
obtain on this must be sought in the mode of election (§§ 30---40). The Folkething's
members are elected directly by those in the population
entitled to vote ``in the proportion of one member
to 16,000 inhabitants.'' The election of the Landsthing's
members, on the other hand, is more restricted and indirect
(§ 34); the whole number is only 66, of whom
the King appoints 12, etc.

That the Constitution by this double mode of election
has wished to secure for itself as many-sided and complete a
representation of the whole popular will as possible is
beyond all doubt. It has for this purpose conceived
the popular will from two different, yet inseparable
principal sides, namely as majority will and as a rational will
grounded in the welfare of the whole. It has indeed
assumed that the rational element of the administration of the state would be just as
fully present in the one Thing as in the
other, but at the same time also assumed that the rational will running through
the majority could best be
tested, clarified and grounded when, by a different
mode of election, the different views were secured
a due influence on the course of
the proceedings. By giving the direct elections a relatively
wide extent the law has provided that the whole population
can, so to speak, get to speak to the government and
% --- p. 7 ---
\setcounter{footnote}{0}%
\opage{7}through its elected spokesmen obtain a decisive
voice in the business of legislation. The living people
with its moods, thoughts and opinions, its changing
needs and demands has then each year its
spokesmen in the Rigsdag: that is, generally viewed,
the significance of the Folkething.

But according to the intention of the Constitution the Folkething cannot
be said to represent the people in everything and in every respect; this Thing is,
by virtue of its mode of election, an expression of the popular will, understood most nearly
as the immediate, mobile majority will. The Constitution
proceeds from the view that a people is not always served by
having everything instantly put through that a majority
ruling at a given time demands; the Constitution
proceeds from the view that the electorate standing behind it, though
fundamentally it always wills what is right, yet does not always have
the full, clear, sufficient insight into what in
every matter is the right thing, and for that reason, when minds
are set in commotion and party interests are in play,
can let itself be carried away into zeal for something which, more closely
considered, is to the harm of country and realm. The Constitution
proceeds from the view that an aroused majority, just because it is
lively, enterprising and aroused, can be misled into taking
steps which the people afterwards, when the consequences have shown
themselves, must bitterly regret.

In consideration of this the Constitution, as a counterweight
to the stirring, forward-driving element in the Folkething,
has by a mode of election of its own, an indirect one, established another
Thing, the Landsthing, whose significance is: to supply the
more restraining, controlling element and by its
majority to step in as mediator between the will of the government
and the majority will.

% --- p. 8 ---
\setcounter{footnote}{0}%
\opage{8}Between the two Things stands the royal government,
embodied in a responsible ministry.

Here, then, in this governmental whole there are three centers,
three authorities, three wills; but where is the ruling
unity? Which of the three authorities has
here the original preponderance? Where lies the proper
center of gravity? It is not easy to
say beforehand; for each of the three has its own independent validity.
That the Folkething may by no means be regarded as a subordinate
member is plain from § 48, which says:
``Proposals for the Finance Act and the supplementary appropriations acts
are first dealt with in the Folkething.'' --- But that the Landsthing is in every respect placed on a level with the Folkething
can be seen expressly from every paragraph in which
the special activity of the Things, each for itself, and their
mutual relation is mentioned (§§ 50---54). If the Folkething can
alter the proposal for the Finance Act, then
the Landsthing, if it finds reason for it, can say No
to the alterations made. Otherwise it would
be the Folkething alone that had the right to grant taxes.
The Landsthing would then have to be regarded as
a remnant of ``the advisory Estates;'' but that
the Landsthing, as regards finance, is just as
fully legislative as the Folkething is plain.
If the Landsthing, as regards the Finance Act, had merely
an advisory, not a legislative influence,
then it would also with respect to the other acts of government
stand a whole step lower than
the Folkething; its authority would then thereby cease
to be real legislative authority. But that this
conception stands in obvious conflict with the spirit
and letter of the Constitution is striking when we
% --- p. 9 ---
\setcounter{footnote}{0}%
\opage{9}compare § 2 ``The legislative power rests with the King and
the Rigsdag jointly,'' with § 29: ``The Rigsdag consists
of the Folkething and the Landsthing.''

The division into two has therefore, when we look at the whole
system of government, become a division into three: the Folkething,
the Landsthing and between the two the King's responsible
ministry. Whether this division into three is now in every respect
suited to its purpose, or whether perhaps for the sake of unity
and concord a change ought to take place here,
is a question by itself. The Constitution, once
revised, can be revised again; it does not
give itself out as infallible. It says (§ 95): ``Proposals
for alteration in or addition to the present Constitution
may be put forward in an ordinary as well as in an extraordinary
Rigsdag.'' But --- this stands firm --- as
long as the Constitution is, it must apply unconditionally; it
is for the time being our highest yardstick for what is
lawful and what is unlawful; we have no other.
To put upon the Constitution another meaning than that
which demonstrably corresponds to its spirit and letter is unpopular,
for it is dishonest. To make use of the sound of words
and give the appearance that the Folkething in everything
represents the whole complete popular will is
unpopular, for it is dishonest. To tear a single sentence
out of its context and twist it to such
a meaning that the Landsthing, as regards the
final adoption of the Finance Act, should not have equal
authority with the Folkething is unpopular, for it is
dishonest. And one thing more: to wish to place the responsible
ministry in so one-sided a dependence on a
voting majority that the royal will
% --- p. 10 ---
\setcounter{footnote}{0}%
\opage{10}dwindles away into an empty form, a meaningless symbol, is
also, as we shall now see, unpopular and dishonest.

What does the Constitution teach us about the King's personal
position in the system of government and about the power and significance of
the royal will? ``The King is not responsible, his person
is sacred and inviolable. The ministers are responsible
for the conduct of the government'' (§ 12). When one now
infers from this that the King's name ought not to be improperly
mixed into the strife of the parties, and that the King
exercises his acts of government only through responsible ministers,
etc., then that is quite in
order. But if one goes further and infers from this that
the King's personal will in reality has no influence
on the course of the administration, then that is as
contrary to the spirit and letter of the Constitution as
it well can be.

By personal will we must here, mark you,
by no means think of private wishes or individual
inclinations. A truly personal will always
aims at something essential, generally valid in a personal sense.
The royal will aims personally at
observing, maintaining and carrying through the Revised Constitution
of the Kingdom of Denmark with all provisions flowing from it:
this the King personally wills; this and
only this can he as King personally will. ``Before
the King enters upon the government, he gives in writing
in the Council of State the solemn assurance to keep the Constitution
of the Realm inviolably'' (§ 7). On this presupposition
the King becomes personally the highest executor of
all acts of government: without the royal will the
law is dead. The responsibility of the ministers is therefore
twofold: they are responsible to the Rigsdag, and they
% --- p. 11 ---
\setcounter{footnote}{0}%
\opage{11}are responsible to the King. What does the Constitution say?
``The King appoints and dismisses his ministers. He
determines their number and the division of business
among them.'' (§ 13). If a ministry does not think
it can justify to the Rigsdag its remaining in its
place, it is undeniably authorized to resign.
But the King, according to the Constitution, has also on his
side the right to dismiss ministers who do not hit upon
what he, for the good of the realm, regards as the suitable thing,
and to choose a ministry where one can be
found that will by all lawful means work
to carry through His Majesty's personal
conception of the right, lawful course in the administration. ---

What does the Constitution teach about the King's relation to
the Rigsdag? ``Without the consent of the King the Rigsdag cannot
remain together longer than two months.'' (§ 13).
``The King can convene the Rigsdag to extraordinary
meetings.'' (§ 20). ``The King can dissolve
either the whole Rigsdag or one of its divisions.''
(§ 22). ``If the King has not confirmed a bill passed by the Rigsdag
before the assembling of the next Rigsdag,
it is considered as lapsed.'' (§ 24). It is plain
from this that the functions of the royal will are more
than symbolic, that they are acts of government of sweeping
significance. By conceding to the royal will
so essential a superior authority the Constitution has wished
to secure a unity animating and directing the whole government.
From this point of view it is important to give oneself an account of how
the constitutional monarchy differs in principle
on the one hand from absolute power (the
absolute monarchy), on the other hand from
% --- p. 12 ---
\setcounter{footnote}{0}%
\opage{12}majority power (democracy) --- especially with regard to
the popular will.

In the absolute monarchy the popular will is submerged
in the royal will, for the King stands above the law: the absolute
ruler wields the whole arbitrarily. In democracy
the popular will perishes in the majority will:
it is not the people's sound reason, in so far as it
works its way forward constitutionally, that is here the main thing;
no, it is the mood of the population expressed
in the most votes that is exalted to supreme, governing
power and authority. In the constitutional monarchy
the royal will stands firm on the foundation of the law.
The popular will receives from this side its expression in the form of
a unitary will, an indivisible single will. This single will
is not arbitrary, on the contrary! it is in the highest sense
popular. The royal will wills what the people
wills, because it wills solely what the law wills; it cannot
will anything else, for otherwise it would, according to the constitution,
cease to be royal. In the constitutional
monarchy the popular will receives at the same time its right, limited
expression in the form of the majority will: it is the divisible
popular will, represented by division, which with us
comes to real utterance in the Rigsdag. The intermediate link
between the royal will, the indivisible single will,
and the divisible majority will, represented in part in the two Things,
is the ministry with its double responsibility.

So long, now, as the proposals laid before them by the ministry
find such support in both Things that agreement
can be brought about, so long the majority will continues
to be in accord with itself, and the royal will thereby
satisfied finds so far no
occasion for particular intervention. The ministry's
% --- p. 13 ---
\setcounter{footnote}{0}%
\opage{13}responsibility to the Rigsdag is in full harmony with its responsibility
to the King. But in the case that the Rigsdag
becomes divided against itself, the matter takes on a different
appearance. If a proposal wins approval in the Landsthing
but is rejected by the Folkething, it cannot be put
through, for the Folkething has an indisputable right to
offer resistance. If a proposal wins the Folkething's
approval but is rejected by the Landsthing, it cannot
be put through either, for the Landsthing too has an
indisputable right to resistance. Disagreement between the Things
themselves is a disagreement in the popular will, divisible
through its representatives; it is the popular will which,
in the form of majority will, has become divided against
itself.

But a will divided against itself cannot
act. It is in vain that the one Thing wishes to appeal
to the fact that it represents a larger majority than the
other, that it has the whole electorate at its back and the like;
such declarations of mood have no warrant in the Constitution.
The Constitution takes account only of the relation
in which the two Things, each with its lawfully elected members,
stand to one another. In the case
of strife it therefore becomes necessary that the popular will
in its other form, the unitary will, the royal
will, intervenes and takes measures so that
the strife can be settled.

To bring about agreement, the royal will can,
as is well known, choose one of two ways: either let the ministry
resign or dissolve the Rigsdag. If the
first means is chosen, then its proper significance must obviously
lie in this, that a change of system really takes place.
Such a change of system must, of course, be
% --- p. 14 ---
\setcounter{footnote}{0}%
\opage{14}in full accordance with the spirit and letter of the Constitution,
and may for the rest lean more toward the mood in the one
Thing than in the other; but it ought nevertheless, if
it is to be popular, to be raised above one-sided party considerations.
If the second means is chosen, H.M. the King having
the firmly grounded conviction that a change of system would
be to the harm of country and realm, then the royal
will has power, on behalf of the people, to dissolve either
the whole Rigsdag or one of its divisions.

But --- one asks --- if this last means
does not help either, if by repeated new elections
it shows itself; that the opposing party in the one Thing
grows larger and larger, what then? Will it not follow of necessity
that a change of system must take place
to the advantage of the majority and its
electors? We answer: Yes! it will certainly become necessary.
If the whole constitution is not to be burst asunder,
it must eventually come to this, that the electorate gets
its will respected. Should the sole will try to
take the majority will by force, then the Constitution would be
broken: we should then no longer have a constitutional, but
an absolute monarchy. But --- let us not skip over
the intermediate steps. The express difference the Constitution
makes between dissolving ``the whole Rigsdag'' and
dissolving ``one of its divisions'' must surely have
something to signify. The meaning must be this: when
one of its divisions is dissolved, it is only the members
of that division who are subjected to new election;
if, on the other hand, the whole Rigsdag is dissolved, then all the members
(those of the Landsthing as well as of the Folkething)
must be elected anew. To repeat the dissolution of merely the
one division (the Folkething, e.g.) and to put provisional
% --- p. 15 ---
\setcounter{footnote}{0}%
\opage{15}measures in prospect, in order to bring about agreement,
would surely be a partisan, hence unpopular procedure.
If the elections now change --- mark you,
after repeated dissolutions of the whole Rigsdag --- so
that the Landsthing eventually becomes of one mind with the Folkething,
then this is an irrefutable proof that
the Folkething on the whole has the mood of the electorate
on its side. The royal unitary will has then on its
part helped the majority will to lawful agreement with
itself and thereby given it back its lost power to act.
But if the opposite proves to be the case,
if after repeated new elections it appears that a
majority in the Landsthing --- no matter whether it is larger
or smaller --- continues steadily to stand stiffly against
a majority in the Folkething, then the leaders of the Folkething have
absolutely no right to appeal to the concord of
the electorate as a whole; for the electors to the Landsthing
--- not forgetting H.M. the King --- belong, according to
the unambiguous provision of the Constitution, to the electorate,
just as fully as the electors to the Folkething.

That resistance, resistance and resistance cannot, if
things are to go lawfully, bring about
agreement is obvious in itself. We can,
without unwarranted interference in practical politics, cast
a glance at the position of the parties. From the one side we hear
declarations which run roughly as follows: ``It must
at last become manifest that what is being fought over here is more
than this or that. We are fighting against a system of government
which no longer suits the times. The persons
change, but the worn-out system keeps coming back:
a little change in expression and bearing, but always the old
face! What we want, what we are striving for, is something
% --- p. 16 ---
\setcounter{footnote}{0}%
\opage{16}new and distinctively fresh, that is working its way forward in
the people. Here the strife is not about a few millions alone,
not about the distribution of taxes alone, not about this or
that means of defense alone: the strife is of direction against
direction, of system against system, of spirit against spirit.
What Luther said to his opponents at the negotiations
in Marburg, that we men of the ordinary
people say with a good conscience to our opponents in the
Rigsdag: ``You have another spirit than we.'' --- From the
other side the explanation runs thus: ``Here there is not at all
a strife of system against system, but the question
is of system against party. The question is whether a
majority party, out of misunderstood zeal for what it itself
stamps as popular, shall have the right to let
party spirit loose, carry the mood away with it, cast scorn
on the judgment of experts and appeal to the multitude. Here
resistance is made for the sake of resistance alone. By
throwing the most dissimilar matters together in a heap and tangling the conditions
into a hard knot one makes any
government impossible. Is it for the good of the country, is it
popular, is it reasonable, is it right?''

As the knot is tightened and the tension grows,
one approaches the critical turning point at which provisional
measures may become necessary, namely
a provisional Finance Act. The possibility of such a Finance Act
has been boldly rejected, and people have comforted themselves with the words of the Constitution
(§ 49): ``Before the Finance Act is adopted,
taxes may not be levied. No expenditure may be made
that has not a warrant in the Finance Act or in a supplementary appropriations act.''
Is the issuing of a provisional Finance Act
then straightforwardly lawful? We answer: No! But to appeal to
the lawful, when one has oneself helped to make the
% --- p. 17 ---
\setcounter{footnote}{0}%
\opage{17}lawful impossible, is a self-contradiction that contains no comfort.
Nor is it straightforwardly unlawful. ``In
particularly urgent cases the King may, when
the Rigsdag is not assembled, issue provisional laws.''
(§ 25). The whole emphasis falls on the ``particularly
urgent.'' But whether a convening of
the Rigsdag becomes impossible because the enemy is in the country,
or the convened Rigsdag, on account of the resistance of the parties % print: sammemkaldte (misprint)
to one another, makes itself impossible, the
case is equally urgent. Instead of comforting
oneself with the worn-out phrase: there is no danger, it would surely
be better to look somewhat more closely at § 95, so that the way
along which the majority will, in full consciousness of its
legally limited right, advances to full accord
with the royal unitary will may come
to lie open and clear before the whole population.

It is the foundations of our constitution that we are about
to deepen. There are banks and reefs in these waters:
we must sound, set out sea marks and secure
the sailing. A few gusts do not frighten us; even
if a gale with breakers comes, our ship
will not run aground: abaft the mainmast stands sound
reason, and prudence sits at the helm. As yet
the speed is good, and the danger only a cloud on the horizon; as yet
nothing unlawful has happened; as yet the testimony stands firm,
the honorable testimony of history, that the Danish people,
even in the moment of tension, is a calm, sober, circumspect on all sides
people.

% --- p. 18 ---
\setcounter{footnote}{0}%
\lecture{II.} {Popular Self-Government: People and King.}
\opage{18}A vigorous, sound and many-sided development of popular life is
conditional on popular self-government. When a people has
advanced so far that it has raised itself to the stage of freedom
and has become of age, it can no longer tolerate any
external master or guardian over it. The
people of age wants to have its will; it wants to govern
itself, rule itself, and take the consequences of what it itself
has resolved.

This principle strikes men of the old
school as something highly offensive, almost horrifying.
``What is to become of it,'' they say, ``when a people becomes
so self-willed that it will not let itself be governed:
insolence, unruliness, licentiousness must surely gain the upper hand
day by day.'' --- One readily grasps that a misunderstanding
lies hidden here. It lies in the concept of
governance that there must be some who govern, some who
are governed, and a law according to which they are governed; but it lies
by no means in the concept of governance that the governing element and that
which is governed must be held apart by an insurmountable
barrier. If the latter were unconditionally necessary,
then all self-government would be impossible: no being, whether
in heaven or on earth, would then be in a position
% --- p. 19 ---
\setcounter{footnote}{0}%
\opage{19}to govern itself. It is a superficial notion
that those who govern are free and can do what they
will, while those who must let themselves be governed never get
their will, but are constantly under compulsion. The government
--- it is said --- is free, the people unfree; the authorities
are free, the subjects unfree; the guardian is
free, the ward unfree. If this relation stood immutably
fixed, then there could undeniably never be any talk
of a free people; but then a ward could surely never
become of age either. If the ward is to become of age,
then the guardian must step down. When can the guardian
step down? We answer: when the governing element, which
hitherto was in him, thus outside the ward, is placed
within the ward himself. The ward then no longer has
the guardianship outside him, but within him: the guardianship
does not cease, but it is transferred. Precisely
by having taken the guardian into his inner being has the ward
become of age. He is of age to the same degree
as he has, spiritually understood, eaten, consumed and digested
his guardian.

The same holds of an entire people. It begins
in minority by having the governance, the government,
outside itself as something alien, an external superior.
The transition to self-government does not come about by
a people casting off the government, much as a horse
throws its rider; no, but by its inwardly taking up
the governing element into itself and moving the seat
of the government to its own center. The guardianship,
the external authority standing outside the people, ceases
when the people has become of age, when it has, in other
words, eaten, consumed and digested its guardian.


% --- p. 20 ---
\setcounter{footnote}{0}%
\opage{20}The essential likeness between self-government in the
single individual and in the people has, however, a real
unlikeness that we must not overlook. When the individual
sets about governing himself, he learns at once that in
his nature there is a lower element which is to be brought
into obedience under a higher. There are sensuous inclinations,
impulses and passions that are to be
mastered; only the rational will, guided by insight, and in
a still higher sense the will of the spirit, ought to rule. But
a human being's will is not moved solely by regard to
what is rational; there is a natural will in flesh and blood,
an impulsive will, which branches out into many inclinations:
the head wants one way, the heart another, self-love
one way, the sense of duty another. As the understanding takes
sides now with one, now with another of these
inclinations, party struggles arise, often stormy
election struggles, in the inner being of man; it is then to be shown
how far there is a ruling unity of will or not.
Without a higher unity of will no character
can form. The individual with the many wills becomes
at odds with himself, is swept along now by one mood,
now by another, grows giddy, wavers and reels, and
ends by losing his way in himself. If a character is to
form, there must be a superior, governing
will which is master of the discordant, a royal
will which is able to reconcile the mutually quarrelsome
petty wills. There is in every thoroughly formed
character a constitutional-monarchical constitution.
The governing sole will cannot by an arbitrary
word of command reduce to silence any single one of the
expressions of the impulsive will, nor does it bring order
and peace by letting the divided wills wear down each
% --- p. 21 ---
\setcounter{footnote}{0}%
\opage{21}other in continued strife, and keep itself behind the majority
of the willing inclinations as an empty symbol, a shadow
image of the law. The will of character is the living law;
it enters into the conditions of the psychical-bodily constitution,
it permeates, purifies, limits and
reconciles the mutually conflicting impulses: the more
constitutional the ruling will of unity is, the
more royal it is. All this takes place in an inward
manner, and only when this really takes place has the human being
himself in his power, and is self-governing.

If we apply this to popular self-government, then
the essence is certainly the same; but its outward
appearance is different. Here in society there is
not merely a multiplicity of directions of the will; but
here there are as many wills as individuals. The people is
in the whole population; the will of the people, seen from this
side, is split and divided: the will of the multitude is an assemblage,
a sum of mutually different wills, a swarm of wills,
which, in order to form a true popular will, must
come together into unity. Now when the governing will,
in order to bring about agreement and unity, steps
up externally against the many self-willed wills,
an ominous and threatening tension arises % print: ag (misprint)
between the will of unity and the will of the multitude. So
long as the will of the multitude, more precisely determined as majority will,
by appeal to the power of the majority sets itself
sharply against the sole power of any ruling will whatsoever,
society is at odds with itself, constantly
hovering between tyranny and anarchy. Under
such circumstances there can naturally be
no talk of rational government, still less of
popular self-government. For the sake of government there must
% --- p. 22 ---
\setcounter{footnote}{0}%
\opage{22}be brought about a rational interaction, determined by law,
between the will of unity and the will of the multitude; there must, in
other words, be given a constitution.

What it has cost to get a tolerably rational
constitution worked out, history teaches. Early
the necessity of a governing will of unity was felt.
Already Homer sings: ``Majority rule is not
good; let only One be ruler.'' But early too
men had to learn the consequences of the inhuman,
or rather, all too human, excesses of absolutism:
the peoples have sighed under a heavy yoke, the will of the multitude
has, after wild revolts, constantly sunk back into impotence.


In the Middle Ages the constitution of society was aristocratic.
Only in the great noble families did a majority will stir
in favor of freedom, mark well: the
freedom of the nobles; of freedom for the common man there was
naturally no talk. With what zeal the nobility strove
to make royal power a semblance and to transform the
royal will into a signing symbol of the authority of
the Council of the Realm, of this the period of the royal charters
in the history of our fatherland --- from Christopher II to
Frederik III --- may furnish sufficient testimony.

By reaching out a hand to the burgher estate and giving itself
the appearance of wanting to help up the whole people, the
fettered royal power shook the chains off itself; but what
happened then? Absolutism cowed the nobility and taught it to
bend before the throne. But the bending only made pride
supple, the nobility kept itself aloft as a court nobility,
and the common people sank still deeper. The general will
was absorbed into the state, and the state was absorbed into the royal will.
It was not vain boasting, but actual truth,
% --- p. 23 ---
\setcounter{footnote}{0}%
\opage{23}when Louis XIV said of himself: I am the state
(l'état c'est moi). The people had now got one guardian
instead of many; but the guardian stood outside
the people, he was exalted over the people with divine authority,
and let it feel that it was in One and All
under age. So it went until the Revolution. With the Revolution
(1789) an element broke through which in
our days plays an important role, namely the democratic.
In order to break down absolutism it became necessary
that democracy should overthrow the throne, level the ground
and pave the way for further progress. Then democracy got
the current with it. Expressions such as: will of the multitude,
majority will, have gained such currency in
public opinion that they are --- to the irreparable harm of
the popular cause --- without further ado accepted as synonymous
with the popular will. It is a lopsidedness of outlook
which must gradually be corrected. The strength of democracy is
to assert the majority will, therein it has its right
and significance. But in its strength it also has its
weakness. Building on the foundation of the majority will,
democracy must unceasingly fight for something unattainable, for
it can never reach the unity of the ruling will;
the majority shifts, party rises against party.
The unity that is sought in the connection of the whole determined by law
is weak and wavering: unity of connection
without unity of will, without thoroughgoing self-unity,
lacks sufficient ground. Resting on a democratic
foundation the building of the state is loose in its joints, constantly
exposed to hurricane gusts, when the forces of the majority stir
in the storm of the party elements. Think of America's democratic
party struggles and the election storms that rage at
every change of president; think of France's
% --- p. 24 ---
\setcounter{footnote}{0}%
\opage{24}entangled condition; old and new dynasties intrigue
against one another, while the communistic republic,
``the red specter,'' threatens with the consistent
carrying through of democracy. The democratic element has
been, is still and will always be an indispensable factor
in the free development of society; but democracy is no
complete principle. If democracy is exalted to the sole
valid principle, then it becomes in its essence just as
one-sided, for the development of popular life just as ruinous,
as absolute monarchy. In order to
clear a way for ourselves to a rational insight into popular self-government,
let us try to shed closer light on it.

There are surely many in our time who find it difficult
to reconcile themselves to the thought that constitutional
monarchy should signify a whole, in itself
coherent system of government; on the contrary it must
appear to them as if it were a kind of patchwork, a compromise
practically extorted by the circumstances of the time, a piecing together
of two mutually irreconcilable systems, trimmed
to fit one another, namely absolute
monarchy and pure democracy. One proceeds here
from the wholly mistaken presupposition that absolutism
and democracy, each by itself, are whole principles. That is
a prejudice drawn from a misunderstanding of history. Historically
considered, each of the two systems undeniably gives
the impression of being a complete whole; but the historical
upheavals also give the impression that society
has been able to settle into repose just as little in the
one system as in the other. There is in the more recent history of Europe
nothing that has to such a degree made difficult a
rational use of constitutional monarchy and
hindered the development of popular life as the false
% --- p. 25 ---
\setcounter{footnote}{0}%
\opage{25}conception that this monarchy is supposed to be a mixture
of irreconcilable components. Even the fusing together
of designations such as the ``limited'' and the
``constitutional'' monarchy contributes its share to nourishing
the misunderstanding.

When the Revolution had ended with Napoleon's fall,
the notion of the old patriarchal
absolutism, kingship by the grace of God, came forth again with
renewed strength. The sole rulers of Russia, Austria and Prussia
concluded --- without the co-signature of their ministers ---
in Paris the \emph{so-called Holy Alliance}
(26 September 1815). In this document it is stated that
the three contracting potentates, in accordance
with the words of Scripture, will regard one another as brothers,
help one another in encountered cases of need, and lead their
subjects and armies, whose heads of families Providence
had set them to be, in the same spirit of brotherhood;
they acknowledge that ``the Christian nation, to which they
and their peoples belong, in reality has no other
sovereign than him to whom alone power belongs,
namely God and our divine Redeemer, Jesus
Christ, the Word of the Most High, the Word of Life.'' To the
light-blue gaze of a trusting-hearted piety, so
Christian-paternal a conduct on the part of the mighty of the earth must undeniably
have looked like a sign from heaven, a promise
of peace and felicity for the peoples of Europe
after long and hard sufferings. It might indeed seem
as if the self-ruling sole will had, by an elevated
temperature of religiosity and love, at last become
so thoroughly warmed through and expanded that it was now really able
to set the majority will at rest and to nurture the freedom of the people
at its bosom. But the statesmen smiled; those
% --- p. 26 ---
\setcounter{footnote}{0}%
\opage{26}who knew a little of the course of the world and remembered that
the treaty on the partition of Poland begins precisely: ``\emph{in the name of the
Holy Trinity}'' felt themselves seized by
uncanny forebodings and already glimpsed the claws by
which one knows the lion.

If the freedom of the people was not to perish, then it
was time for the democrats to take the field with
the majority will. It was a course of action that demanded
endurance and courage. Every attempt either to
obtain or even to preserve such rights (franchise,
the right to grant taxes, freedom of the press, freedom of assembly, etc.)
as are necessary if a people is to be able to move
freely, every attempt in a liberal direction was regarded by the
absolutist governments as a brazen assault
on the Holy of Holies. To expose abuses was
impious wickedness, to turn to the multitude was a
political crime. The fanatical adherents of absolutism
saw in every freer movement a new head shooting
forth on the hydra of revolution, and hastened to hack
it off. Here is a turning point where democratic civic spirit,
courage and self-sacrifice show themselves in their most favorable
light. The Revolution had --- fortunately for the cause
of humanity --- given democratism such a mighty upswing
that the free spirit no longer let itself be cowed. In the period
from 1815 to 1830 it is the leaders of democracy who step
into the foreground and must draw the load. Their intrepidity,
zeal and endurance have --- despite all
demagogic excesses --- a rightful claim on the gratitude and
appreciation of posterity.

From this one will be able to see how it has come
about that the two half-principles, the monarchical with the will of unity
and the democratic with the majority will, have
% --- p. 27 ---
\setcounter{footnote}{0}%
\opage{27}given themselves out, each for itself, as a whole, complete principle
and have appeared not merely as temporary opponents,
but as born enemies. From this it will be understood
that the concessions such two powers under continued
conflicts can wring from one another are by no means suited
to form a system of government with true organizing
unity. The constitutional monarchy, the monarchy
surrounded and limited by democratic (republican)
institutions, must, when it is seen from this side,
necessarily suffer from the ills which poisonous substances
incompatible with the organism cause in an otherwise vigorous
body. What convulsions may arise from this
is historically illustrated in the disputes under Louis
XVIII and Charles X over the application and interpretation of
the French Charter.

Nothing is more ruinous than to build on a
Constitution that is both to hold and yet not to hold.
Conceived as a piecing together of two hostile principles,
the constitutional monarchy is precisely suited
to awaken, nourish and maintain an incessant discord
between popular representation and government. When both
parties are at bottom dissatisfied with the constitution, they can
by resistance, intrigues and excesses, each from its own
position, harass one another endlessly.

If it is the royal will which in absolutist selfishness
feels the urge to extend its power, the encroachment from
above will never lack means. Here it does not help to
rely on ministerial responsibility. What cannot
be done by force can be done by cunning; if one ministry does not dare
to expose itself to a sentence of the Court of Impeachment, well, then another
dares: a Charles X will always find a Polignac. There
is among men of government who do not take it strictly with
% --- p. 28 ---
\setcounter{footnote}{0}%
\opage{28}the moral side of the means a peculiar shrewdness in sounding out
moods. One has in all secrecy a party in
reserve, one makes masked movements which confuse
the opponents, and with proud contempt for
public opinion one keeps to the one decisive question:
\emph{how much can one demand of a people before it
rises in revolt}? True, the calculation may go wrong, but
one must venture: the attempt may after all succeed. What
failed for a Charles X (1830) turned out happily
in the coup d'état (1851) for a Napoleon III.

If it is, conversely, a democratically minded popular representation,
indeed even if it is only a section of it,
that is dissatisfied with the constitution, then the encroachment
from below will likewise not lack means for cunning exploitation.
First and foremost it will then be a matter of
using the shortsightedness of the multitude, or at any rate of the majority,
to get the Constitution, held at a distance, set in
so skewed a light that it looks to the life as
if it spoke unconditionally for the democratic view.
Detached sentences, slogans and turns of phrase are well
suited to agitating the population. The common
people believe sincerely in such explanations --- especially when
they do not know, and are warned in the most emphatic manner
against making acquaintance with, other explanations,
since these are known in advance to be a web of
semblance and deception --- and are then easily ensnared in the delusion
that their rights are being violated, that the democratic leaders,
by no means from greed for power, but in true
self-denial and with much self-sacrifice are fighting for a
true, a whole and complete principle for the benefit and
good of all. Next --- when the agitation is in full
swing --- the thing is to intimidate the government and
% --- p. 29 ---
\setcounter{footnote}{0}%
\opage{29}to overthrow one ministry after another. The latter
means is as popular as the first. History
teaches that the common people, so long as it has not yet
become of age, gladly rejoices and claps its hands
when a government falls. There is a series of
examples that this has happened even in such
cases where the fallen government was the actual
defender of the people's cause, and it was the enemies of the people who, through
the upheaval, came to the helm. If a government under
such circumstances lets itself be intimidated by the
impudence of majority leaders, it may quite certainly end in
an overthrow of the constitution. If a government on the other hand,
by holding unshakably to the foundation of the law,
transforms the democratic assaults into rodomontades\footnote{Great boastings.}
and blind alarm, it may end with the general
sense for great interests gradually being dulled, and the tumult
giving way to a drowsy calm. The latter is, in the popular
sense, just as sad as the former.

We have so far assumed that it was only one of the parties
who by encroachment wished to violate the constitution;
but such a case belongs to the rarer ones. It is
far more usual that encroachment from the one side
provokes encroachment from the other. We have an example
of this in the history of Sweden in the period from the
death of Charles XII to the Revolution under Gustavus III. ``The
death of Charles XII and the election of his sister Ulrika Eleonora as queen
(1719),'' --- it says\footnote{Verdenshistorie efter Cæsare Cantù, frit bearbeidet af E. Holm og Ph. Weilbach. P. 498 ff.} --- ``had made Sweden
from an absolute state into an almost republican
% --- p. 30 ---
\setcounter{footnote}{0}%
\opage{30}realm. Ulrika Eleonora gave up the government already the year after
her accession to the throne, having it handed over
to her husband Frederick of Hesse, a prince who by no means
lacked abilities, but was superficial, pleasure-seeking
and always stood as a stranger in his realm. Hatred of
absolutism and displeasure at the devious ways by which Frederick
and his queen strove to get as much power
as possible, brought forth at several Rigsdags (1720 and 1723)
a series of such constitutional provisions that one
would have difficulty pointing to any other kingship
in which the King has had less to say or been placed
in a manner more oppressive to his personal
dignity\dots{} A strongly agitated political life developed
in Sweden under this republican constitution.
A stirring and a practice in critical
discussion was created, which spread to other domains
to the benefit of the development of the feeling of independence and
of free movement. Knowledge was gained of the parliamentary
forms, and through the numerous conflicts
a series of experiences was made, which have turned to the benefit of
political life later in Sweden. The thought that the people
had a right to set up legally fixed barriers
against royal power took root, never again to vanish.
But the one-sidedness with which the weight was laid on the
Estates, in connection with other unfortunate circumstances, caused
such strong shadow sides gradually to come forth
in the ``era of freedom,'' that one can understand the displeasure
with which most Swedes look back on it.
--- When Frederick of Hesse was dead, Adolf Frederick, according to the established
order of succession, ascended the throne (1751).
The indifference and sluggishness with which his predecessor
had almost let royal power fall into oblivion
% --- p. 31 ---
\setcounter{footnote}{0}%
\opage{31}was now replaced by attempts to extend the King's influence,
the result of which was his still deeper fall. Adolf
Frederick's mild, good-natured character and strongly limited
mental abilities seemed to make him well suited to content himself
with the modest latitude that the constitution granted
him; but his queen, the sister of the Prussian king Frederick
II, Louisa Ulrika, who had complete
mastery over him, was as ambitious
as gifted.'' --- To what degree the queen's plan for
a revolution nevertheless failed, is seen from what follows. ``The
imprudence of a couple of subordinate participants caused
the plan to become known, and the punishment followed hard upon the
discovery. The King, who had known nothing of the plan,
himself in a letter to the Council broke the staff over the
men who had risked their lives for him, but did not thereby save
himself and the queen from new humiliations.
While eight of the participants in the conspiracy were
executed, the clergy gave the heavily compromised
queen a sharp reprimand, in which it let
her hear that she had set aside her duty to God,
the King and the realm of Sweden. The King himself had to see it
declared in a letter from the Estates that he had broken
his oath, so that the Rigsdag, if it wished, had a right to
declare that bond to be loosed which tied the subjects
to him.'' --- What did not succeed for Adolf
Frederick and Louisa Ulrika, succeeded for their son, Gustavus
III, who after his father's death (1771) ascended the throne.
When all more lenient attempts to smooth out the passionate quarrels
of the parties (the Hats and the Caps) failed,
his resolution matured to venture a revolution.
Without waiting for the arrival of Sprengtporten and the Finnish
soldiers, he understood how ``by his brisk, inspiring
% --- p. 32 ---
\setcounter{footnote}{0}%
\opage{32}conduct to win over to himself the troops that were in Stockholm.
The burghers joined him, while the Council of the Realm
and the secret committee were at a loss, so that the Revolution
was scarcely begun before it ended victorious.
After all resistance had fallen silent, the
King delivered a thundering speech to the assembled Estates on
the misfortunes which the era of freedom had brought upon the country.
Thereupon he had a new constitution read aloud, and when he
from the throne, with Gustavus Adolphus's silver hammer in his hand,
asked the Estates whether they consented to it, a thrice
repeated Yes rang out in answer. The same Rigsdag
which a few days before had shown the most reckless
arrogance, now did not dare the faintest resistance
to protect that freedom of the people of which it,
as bearer and defender, had been wont to boast of acting.
Outside Stockholm the people joined with joy
the new order of things, so that the constitutional edifice of the era of freedom
collapsed like a house of cards, which
topples at the first breath.''

From a popular point of view the constitutional
monarchy puts itself in an entirely different light. The
monarchical is not a clipped remnant of the absolutism of absolute
monarchy, any more than the constitutional is a
product of a half-tamed democratism; but both
are sides and products of the very essence of the popular will.
The will of unity, represented and brought to life in the person
of the King, is just as organically necessary a member of a
popularly free constitution as the majority will represented
by the delegates of the population. The Constitution must therefore
be framed in such a way that the popular will in the
two principal forms can come forth powerfully and keep itself in
constant organic interaction with itself, determined by law.
% --- p. 33 ---
\setcounter{footnote}{0}%
\opage{33}There is danger afoot when the repeated majority will, expressed in lawful,
unambiguous form, is violated; but
there is also danger for the cause of the people when the lawful
form for the decisions of the will of unity is not respected.

So long as a people has not yet become of age,
it will constantly live on in the delusion that King,
government, officials, intelligence and whatever belongs under that,
together make up an alien power, which constantly,
if not always openly, yet secretly, aims
at cowing and oppressing the people as much as possible;
that is to say: the people in its state of minority has not
got further than to conceive and feel itself as the common people.
This minority the democrat knows excellently how to exploit,
when he steps forth as demagogue. What is a
man of the people? It is a man of the common people; everyone
who, whether by birth or by fortune or by intelligence,
raises himself a little above the level of the common people already stands
outside the people and is a suspicious person. The demagogic
road to popular self-government is clearly marked
hereby: by striking down the guardian the
one under age is to govern himself.

The more a people works its way forward to majority
and the maturity of manhood, the more clearly it will see that
the governing power, which must necessarily stand above the level of the common people,
does not fall outside, but within the domain of its
own will. A people of age will no more
overlook the right of the common people to effective, lawful
participation in the whole development of society than it will in
One and All reduce itself to common people and make the government of the state
a plain and simple affair of the common people. A people of age
will see to it that the government and its officials do not
act arbitrarily in a bureaucratic spirit
% --- p. 34 ---
\setcounter{footnote}{0}%
\opage{34}opposed to the popular will; it will make sure that the state does not
appoint or pay its officials in a manner
out of proportion to the measure of spiritual forces
they apply in the service of society; but when the
proper safeguards in this respect have been applied,
then the people will also know how to hold the governmental authority
and its servants in esteem; for the officials of the
state are indeed the people's own officials, organs of the popular will. A people that knows what it is doing
will just as surely find it insulting, indeed outrageous,
if anyone advises it to weaken the organs of its own will,
as a reasonable man will find it outrageous
if anyone is so foolish as to advise him to weaken
the muscles and sever the nerves by which he
alone is able to move his limbs. For
a people of age it stands clear that the executive
power too ought to have a free domain, that the authority of those who govern
is the people's own authority, the intelligence of the teachers
its own intelligence, etc.

The constitutional cooperation between the will of unity
and the majority will must be so sufficiently
determined in the Constitution that every encroachment from
the one side as well as from the other can be combated and repelled
in a constitutional way. This is, in brief,
the meaning of \emph{popular self-government}.

The King's will and the people's will are not two, but one
will; the majesty of the people and the majesty of the King are not
two, but one majesty. In the royal majesty the
people honors its own sovereignty; it beholds in the King's
person its own personal unity. The King's freedom
and the people's freedom do not limit one another; it is
% --- p. 35 ---
\setcounter{footnote}{0}%
\opage{35}one and the same popular constitutional freedom. That the King
stands over the people means that the law stands over the people.
The King's will is the execution of the law: the King is the
living law; in this sense we understand the words: ``\emph{a free
King over a free people}.''

% --- p. 36 ---
\setcounter{footnote}{0}%
\lecture{III.} {Political Experiences: Parliamentarism.}
\opage{36}Experience, it is said, is a good schoolmaster, and that
is true; but the schoolmaster is sometimes rather dear,
according to the old saying: through harm one becomes wise, but seldom
rich. If one at least always became wise, the harm
might be borne; but it is not said that one
always becomes wise, one can unfortunately also become stupid
from sheer experience.

It is not so straightforward a matter to gain experiences,
as many imagine, that is to say: if the experiences
are to be worth anything. To overcome difficulties,
one must know them. As an introduction to a closer
consideration of the difficulties that stand in the way
of those who devote themselves to political experiences, we might
begin with some experiences drawn from other
circles.

Natural science is a science of experience; there are facts
enough here, and enough to experience. But what has it
not cost before research got so far
that it could determine the right procedure and
secure itself against false inferences. It is an old
experience that sometimes shortly before sunrise
% --- p. 37 ---
\setcounter{footnote}{0}%
\opage{37}a bright star shows itself in the eastern sky, i.e.
the Morning Star, and at other times immediately after
sunset likewise a bright star in the western
sky, namely the Evening Star. The opinion lay
close at hand that these must be two different stars;
but although it looks like a real experience, it is
nevertheless only a false inference. Antiquity was surprised
when one of its sages discovered that the Morning Star and the
Evening Star were one and the same star. It is an
old experience that stones fall and that smoke rises
into the air. The Greek philosopher Aristotle built upon
such experiences the doctrine that some bodies by nature
moved downward, others upward, and that
an essential distinction had to be made between bodies
with and without weight. But physics teaches us that the notion
of bodies without weight rests on a
false experience. Experiences in the sensible world
are not gained merely by hearing, seeing and feeling; it is necessary
to be able to think, judge and infer.

If in the realm of nature, where after all everything proceeds by necessary
laws, it is difficult to find one's way in the
multiplicity of experiences, how difficult must it not be
in the world of freedom, in human life, where a choice takes
place, and where the will's choice does not always follow the
commands of reason, but prefers to let itself be guided by impulses
and incalculable inclinations. Here
it depends not only on what is experienced, but
first and last on how it is experienced. It is a well-known
matter that two people can live through the same event
and yet get the most widely different things out of it.
Inferences from experience are often in a high degree hasty and
crooked, although the person concerned, who imagines that he has
% --- p. 38 ---
\setcounter{footnote}{0}%
\opage{38}facts on his side, himself regards them as unfailing.
It is an experience that men can be deceived by
women, and women by men. The man who is so
unhappy as to be deceived once, twice, indeed
perhaps three times, receives, if it really goes to his
heart, experience's certificate that three times is all
times, and becomes a woman-hater. Women, whose
finer nerves do not tolerate any repetition, can have
enough in a single experience to become man-haters.
Indeed, if one may rely on the poet --- and the poet
must surely understand such things --- there are said even to be
female individualities with so finely discerning a
sense that, precisely because they have not had the luck themselves
to have experiences, they know from experience that all men
are deceivers. Of the experienced aunt who had undertaken
to watch over ``the Rose of the Alhambra,'' Washington Irving
remarks: ``The ever-watchful Fredegunde was
the most prudent old creature one could imagine.
With regard to `the hostile sex,' as she called
it, she cherished a respectable horror and distrust,
which had gradually increased with the prolongation of her unmarried state. The good lady
herself had indeed never suffered from the sex's intrigues, since Nature in
her face had given her a protection against
all pursuit; but women who have least
cause to fear on their own account are
precisely the most ready to watch over their more
tempting fellow creatures.''

As regards the upbringing of children, one is not
seldom witness to how mutually discordant experiences can
be. The one party will have it from experience that the child
must be kept under strict discipline if it is not to become unruly
% --- p. 39 ---
\setcounter{footnote}{0}%
\opage{39}and wilful, the other party can prove from experience
that the dear child only becomes defiant when one takes
it in the bad way, but is always well-behaved when one yields
a little to it and takes it in the good way.

When the talk is of popular enlightenment and learned instruction,
it can sometimes, especially in lively discussions, make
one uneasy to see how experience stands up against
experience. We hold ourselves assured that the people's sound
sense is open to the highest interests of life, that the common
people have a quick, instinctively sure power of apprehension,
a natural keen understanding, which errs only
when it is entangled in unclear moods and led astray by
groundless notions. It is, in our view,
a pressing necessity that the common people,
the great mass of the ordinary population, be affected by the word of light
and truth and get their eyes opened; it is
our hope that the cause will prosper when there are teachers enough
who know what they are about and speak plainly, without flattering
the multitude and fishing for majorities. But can this be straightforwardly
proved by experience? Not quite. There are
thoroughly learned men who think they have experience that
so-called popular enlightenment does harm instead of doing
good. ``The surveys of the unlearned are only overhaulings, which
lack thoroughness. And why? Because they lack
criticism, archaeology, philology and grammar. Loose
hints and scattered gleams of light set minds fermenting
and make the man in the street conceited. Horses that work
hard ought not to be fed with grass: they must have
oats. The oat-kernel of instruction lies in archaeology
and philology; the ``people-awakening'' enlightenment is
a green fodder, which blows up but gives no strength.
The kernel of all that the ordinary people are to learn is briefly
% --- p. 40 ---
\setcounter{footnote}{0}%
\opage{40}and simply this: they are to learn to look up to learning
and show respect for the experts. Nothing is
more pernicious than ``the living word,'' with which one
stirs about in everything. What is this ``living word''?
It is mushy chatter. But mushy chatter is like mushy food;
the living word puffs one up like cabbage and can so easily
cause bloat. The puffed-up, bloated spokesmen
are half-educated rascals, who join in the talk about everything,
but know their way about in nothing.'' --- See, those were now learned
experiences, and it is truly no joke to take up
the cudgels against learned experiences.

Against the learned experiences stand the unlearned. We
respect learning and science, in whatever fields the investigations
may move, be it in the archaeological,
the philological or the geological. But there
are people of a popular awakening who have experienced to such a degree how
hollow, unfruitful and empty learned research is, that they
on the spot can read them the lesson. ``What good does it do,''
they say, ``to dig in strata of earth and rummage in old
documents? Can one dig up life as one digs up
corpses? Collections from the Stone Age, the Bronze Age,
the Iron Age, etc., may well be entertaining enough
for those who let themselves be buried in a museum, but we,
who breathe under God's free heaven and do not feel the need
to be buried alive, what gain have we
from them for our present life? If from the life
our forefathers lived in the Stone Age one is to work out how
we men of the present are to go about
living in the Age of Politics, then more than one
bushel of peas is needed. People gnaw at the many languages and the many
grammars: Chinese, Persian, Tamil,
Malabar, Sanskrit, Prakrit, language upon language,
% --- p. 41 ---
\setcounter{footnote}{0}%
\opage{41}grammar upon grammar! What good does it all do, when
in the end one cannot make living use
of one's own living language? The great junk-heap of antiquity is
a great ``kitchen midden,'' and the endless language-research
an enormous grammar. Let one lay the archaeological
midden upon the grammar, or, if one prefers, the grammar
upon the midden; let one call physics to one's aid
and let an electric current pass through --- from the grammar
to the midden, from the midden to the grammar
--- yet there comes no gleam of life, no spark of
spirit: the dead is dead and remains dead, however long
one may research and dig.'' --- These are unlearned experiences,
and unlearned experiences are by no means easy to combat.

If we take in the Holy as well and reflect upon the religious
experiences, then the strife is still more pervasive.
The many who gathered about Christ, and
were eyewitnesses of his miracles, went away with highly
different experiences. Some experienced that it was in truth
the works of God; but others experienced that it
was blind fanaticism, and the whole an uprising against God
and the Emperor. How much did Pilate experience, when he
asked: \emph{what is truth}? It is told in the Acts of the Apostles
(Chap. 25) that when King Agrippa and
Berenice came to Caesarea to greet Festus,
he was in perplexity and did not know what he should
judge about the Jews' accusation against Paul. The only thing
he had experienced was that they had some disputes
about their worship of God and ``about one Jesus, who is dead,
whom Paul said to be alive.'' That is a characteristic
trait. The whole event of Christ's death and
resurrection stood before Festus as something so insignificant
and transitory that it was unworthy of a Roman to
% --- p. 42 ---
\setcounter{footnote}{0}%
\opage{42}acquaint himself more closely with the matter. So it went for
more than two centuries with the Roman experiences
and Christianity. The party of Christ was from the start
regarded as a mad party of the common people; of such a
party of the common people it was then permitted to believe the most absurd, % print: var. (misprint, stop for comma)
the most foolish thing that could be invented: one had, after all,
experience on one's side.

We now turn to the subject of political experiences.
These experiences are beset with all the
difficulties we have met in what has gone before. They
are to conform to universally valid laws; but the
laws of freedom are so mobile that inclinations and
moods tempt to hasty inferences and can give
us a wholly crooked conception of the real conditions.
To gain true experiences an unbiased mind is required,
one that is able to free itself from all preconceived opinions;
but where personal interests are in play,
preconceived opinions readily settle into place as
unfailing truths, inviolable principles.
To gain true political experiences it is not
only necessary that those concerned should in general
have respect and feeling for truth and right: here
something more is still required. Here self-conquest is required,
in order to put oneself in the opponent's place and to understand
his view from the ground up. Such a self-conquest
is not to everyone's taste. The initiated are just as
tempted to remain in their once-closed circle of ideas
as the uninitiated in theirs. It is so
agreeable to live with nothing but like-minded people and to hear
the surroundings' echo of one's own opinions. From this it can
be understood that mighty parties can fight for a whole
age and with great acuteness spy out
% --- p. 43 ---
\setcounter{footnote}{0}%
\opage{43}each other's wrong sides without managing to see through
each other's nature. We shall consider the difficulties
first from above, next from below.

If we look from above downward, it is surely a matter of course
that a King on his Throne must gain political experiences.
But the experiences, unfortunately, become colored
when they are seen through colored glass. There is an atmosphere
called the court atmosphere: it has a
peculiar light-refracting effect. It is well known that the light
from a star is bent by passing through different
layers of air, so that it looks as if the star stood
somewhat higher than it in reality stands. Such an
appearance rests on what in astronomy
is called refraction. The nearer the stars are to the
horizon, the more the appearance has effect, so that they
seem to stand higher. The astronomer knows that he must
calculate the refraction. For the monarch's observation it will,
on account of the court atmosphere's light-refracting power,
only too easily look as if the lower classes of society stood
higher, in other words: as if matters stood
better with the subjects in his kingdom than they
really do; it is by no means always convenient for him to
calculate the refraction. When Catherine II had conquered
the Crimea, and 30,000 Tartars had been cut down, Potemkin
wished to show his imperial mistress a magnificent
spectacle of splendor and glory. The Empress made a
journey to visit her newly conquered lands. ``Potemkin
assembled a great army on the Dnieper and by the help of skilful
scene-painters he fabricated, so to speak, a
flourishing condition of the country, such as it had never
known. The banks were strewn with villages,
but they were painted on canvas: one saw churches
% --- p. 44 ---
\setcounter{footnote}{0}%
\opage{44}under construction, ships and houses. With the knout
Tartars were driven from far away down to the river banks, so that
these might seem populated; the young grass was
gnawed off by cattle that had been fetched from a circuit of 400 versts.
There would scarcely have been greater expense in
getting something useful established. But Catherine let herself be
deceived in order to be able to deceive Europe with regard to the empire's
strength and her own activity.'' --- What
dazzles a King's eyes is, in rarer cases,
base flatteries; for in that respect experience will surely
teach him to discount. But if there is a circle of
confidential advisers, on whose infallibility he, without
having himself sufficient insight into the course of affairs,
will rely unconditionally, such a circle can become dangerous to him.
When the sky is overcast with a veil of light clouds,
rings and mock suns regularly form about the sun;
so too in the court's vapor-sphere mock suns will
easily form about the royal sun; a confessor, a
first minister or whatever the oracle may be called, will see to it
that His Majesty never sees naked reality,
but images in a mirror, which throws back the glare of his
own light. Even if reality becomes so
pressing that by an upheaval it tears the princely
person out of the magic circle and lets him pass through
the severe school of great trials, it is
nevertheless not thereby given that he gains true experiences.
Let one think, to recall the history of the Stuarts, of
Charles II. What had this Prince not lived through. His
royal father had to lay his head on the scaffold;
he himself had had to go through, in several years of exile,
the humiliations that always meet fallen
greatness. And when he then at last --- after Cromwell's
% --- p. 45 ---
\setcounter{footnote}{0}%
\opage{45}death --- was called back and entered his capital
amid the jubilation of the people, what experiences did he then
bring with him up onto the Throne? ``That every woman's honor was
for sale, and every man's fidelity to be had for money.
The difference was only in the price: some sold themselves for a cheap
bargain, others let themselves be paid dearly.'' Let one read
Macaulay's grand-scale account of Charles II's last
days, and one will see what sort of experiences of life
this King had reaped. Of the Bourbons Napoleon has aptly said: % print: Bourbornerne (misprint)
``they have forgotten nothing and they have
learned nothing.'' They have thus, despite revolution and
empire, only turned about in the circle of their own imaginings,
in old warped experiences.

If we look from below upward, would it then perhaps be
easier for the peoples to gain true political experiences
than it is for the Princes? By no means. If positions of power
in society could be grouped like trees in
the forest, if the course of government, up and down, were as plain
to see as the swaying of the reeds at the edge of a lake or
as the waving of the corn in the field: then sound
understanding would easily make its way and distinguish between right and
wrong. But the struggle for power is an intricate struggle;
no story of seduction is so long-winded as that which
deals with mutual political unfaithfulness between the peoples
and their leaders.

Most often it is probably the leader who abuses the people's
trust and betrays the good cause when unforeseen
turns occur. Yet it does not as a rule happen in such a way
that the person concerned steps forward with the fully conscious
purpose of wanting to deceive, of wanting to ingratiate himself
and win popular favor merely in order himself to come to power.
We must, for the honor of the human race, on the contrary
% --- p. 46 ---
\setcounter{footnote}{0}%
\opage{46}assume that most of those who turn to
the people begin in honest sympathy with the people. But
for all that it may nonetheless end in deception. The
moral self-esteem which is awakened in the leader by the fact
that the multitude joins itself to him is exposed to great temptations.
The celebrated man already feels himself a holder of power, and
this feeling is very tickling; it drives him forward:
he cannot be content with moral influence, he needs
fist-power; he must of necessity command
a party and choose the means as they come. Then
he is irritated by opponents, then he is egged on by worshippers, then
he becomes a hero of the masses, Commander of the Order of the Majority,
and the despot in him awakens. It can then happen in the
most natural way that he who begins with the people's
cause may pass over to serving his own cause and
end by betraying the people.

But the relation can also turn in such a way that
it is the people who betray the leader. There is in the
common people --- so long as they still find themselves on
the stage of minority --- a peculiar alternation between
credulity and distrust. By nature an uncorrupted
people is credulous; that holds as a rule. Anyone who
steps forward and wants to take up its cause easily finds support.
But the blinder the trust and enthusiasm are
with which it follows its leader, the more precipitous again
is the reversal by which it gives him up when its suspicion
awakens. The history of Greece and Rome has
eloquent examples of this. The need for a reorganization of
land-ownership and the strife over the Agrarian Law was for the Romans,
a little over a century before the birth of Christ,
what the questions of pauperism and poor relief
are in our day. The brothers Tiberius and
% --- p. 47 ---
\setcounter{footnote}{0}%
\opage{47}Caius Gracchus, gifted men of noble birth, with
rare virtues and the most honest intentions, exerted every effort
to help the poor impoverished people onto their feet;
but their endeavors were fruitless, they had to see
themselves branded as criminals and pay with their lives.
The nobility with the Senate understood excellently how to lead the mob
by the nose: the ignorant common people were stupid enough to
let themselves be dazzled, cowardly enough to betray their protectors.

A necessary condition for a people to be able to make
progress in political experiences is, among other things, this,
that the discussions of public affairs be conducted
with deliberation, so that the grounds for and against may come
to the knowledge of the public; but that is precisely a condition
that is difficult to fulfill. When parties get
into a passion, one gets no time to think about the matter;
one tumbles about with the personal: everyone is busy
attacking, slandering and blackening his opponent.
Readers who are used to the blackening may well know % print: Læsere. (misprint, stop for comma)
from experience that pressing accusations are not
always to be taken literally; but those who speculate
in excitements also know from experience that a bold
distortion seldom misses its momentary effect.
Democratic America has a knack for applying the
method. Washington took it calmly. Though
he was violently attacked by the Democratic party,
he was nevertheless not angered; he wrote only to its leader,
Jefferson: ``I would never have thought, I will not
say that it was probable, but merely that it was possible,
that, while I was using my best powers to
create a popular policy that was entirely our own, and
to secure the country against the horrors of war, all my actions
should be distorted, represented in a
% --- p. 48 ---
\setcounter{footnote}{0}%
\opage{48}gross and insulting manner, in such exaggerated and
unseemly terms that they could scarcely be applied
even to a Nero, to a notorious bandit, but rather
to an everyday scoundrel. But enough, I have
already gone too far in uttering my feelings.''

If you would gain \emph{true experiences, then master
your passions and use your understanding}. This
Washingtonian principle is but little attractive
to the great mob, and therein lies the knot. It is
not want of knowledge, for that can be remedied,
but want of self-control that confuses
the people's sense and prepares so many obstacles for the popular cause.
Naturally nothing can be objected to the
leaders of a party presenting the matter from their point of view
and seeking to support their conception by grounds as strong,
penetrating and convincing as possible; but
to pass over the opposing party's grounds in silence and to present his utterances
as simple, indecent wretchedness, is
that honest, is that popular? That one wants to hold up a mirror
before the people's eyes is in order; but it is, as is
well known, not all mirrors whose images correspond to the objects.
One has in gardens the so-called globe mirrors.
However well-shaped the strollers may be, when
they come near to such a mirror they undergo
a dreadful change: big noses, fat bellies,
tiny little legs! they look like nothing but nutcrackers.
There are other mirrors that glorify the figures and show
them in supernatural size. Now when the party leaders, the
men who give themselves out as holding up the mirror of truth
before the eyes of the multitude, constantly combine the two kinds of
mirrors in such a way that they themselves come to appear as
heroes, while their opponents stand as
% --- p. 49 ---
\setcounter{footnote}{0}%
\opage{49}nutcrackers, is that guiding the people to political,
true political experiences?

On progress in political experiences rests the influence
of so-called parliamentarism. We say:
the influence; for the origin coincides with the laying of the foundation
of any freer constitution.
Already with the introduction of advisory Estates there follows
a certain kind of parliamentarism. The first thing characteristic
of parliamentarism is that the most important
affairs of government are made the subject of a negotiation
between the government and the popular representation. The popular representation's
influence can pass through different
degrees; but only when the outcome of the negotiations is decided
by a vote in such a way that the majority's vote
at the same time becomes decisive for a ministry's continuing
in existence or not, does parliamentarism come into
full force. For the carrying through of parliamentarism it is
necessary that the view expressed by the majority in the vote
gain an actual and legal influence
on the course of affairs; but it is by no means necessary
that the majority should be gathered in one chamber; it can
well also be divided between two. To stay with
our own Constitution, it is not seen that it makes
any breach in parliamentarism by letting the Folkething
and the Landsthing each have its own majority; for it may
surely be presupposed that a ministry that gets the majority of both
Things against it must resign. Something quite
different is it that provisions are still lacking for
what is lawfully to be undertaken in case of continuing
disagreement between the divided majority. The indefinite
notion of the possibility of provisional measures
is dreadful not only for the reason that in
% --- p. 50 ---
\setcounter{footnote}{0}%
\opage{50}what is to show itself when the law does not suffice
there is always something bogeyman-like; but especially for the
reason that the majority will can neither come into its
full right nor be legally limited when
the Constitution does not definitely and clearly indicate the way
it has to go to reach the goal. As a model
for a consistently carried-through parliamentarism
one is wont to extol the English constitution. But
that the English Lower House can just as little be assumed
to be a pattern for our Folkething as the Upper House
``can be a pattern for our Landsthing, is obvious
to anyone who has even a merely superficial acquaintance
with the English Parliament.

Parliamentarism in England has developed in a peculiar
way out of the particular demands of English conditions
with historical necessity: it has the political
experiences of centuries as its foundation. Brought up to self-government,
the English people has learned to keep strictly
to the letter of the law and to respect the right of the Crown. And yet
the swell can sometimes become dangerous enough when the sea of the people
comes into commotion. The danger is, however, passing. With
the exception of some philanthropic paroxysms, it is as a
rule not great ideas but great projects that rule
the mood. The law has a bulwark in capital; for
to carry out great projects money is required. It is then
also a piece of luck for England that in stirring times it
has a guide as practical as the \textit{Times}, a weathercock
above the Exchange, so politically all-round-looking that it, without
changing its standpoint, can constantly turn with the wind.

On the Continent parliamentarism has in recent
years not exactly reaped laurels. In 1870 it was surely
in good part the revived parliamentarism that,
% --- p. 51 ---
\setcounter{footnote}{0}%
\opage{51}in connection with the moral importunity of opinion,
by inflammatory means tore France into the war with
Germany.

A closer consideration of the nature of parliamentarism
must lead to the recognition that its working
is conditioned by the cooperation of two related, yet dissimilar
factors; the one is the moral, the
other the juridical. The moral factor lies in
the negotiations; for what should the discussion and negotiations
be aiming at, if the intention were not
to have the grounds that speak for and against a pending
matter tested conscientiously and carefully? The
view that is supported by the most convincing grounds
must then be said to have the moral preponderance. The
juridical factor lies in the outcome of the vote. When
a vote has taken place, that view must prevail
which has the most votes on its side; that is
juridical right: no reasoning helps here. If
the two factors now always cooperated in full agreement,
then a consistently carried-through parliamentarism
would undeniably be a political blessing for country
and people. But unfortunately, experience teaches that the opposite
is often the case. The moral factor is sometimes
encumbered with immoral elements.

There are indeed party views that stand so firm that
they do not let themselves be shaken by counter-grounds. Even if
eternal reason were to condescend to take on human shape
and take part in the negotiations, it would
still not be able to make the party uncertain. The negotiations
are conducted only for appearance's sake, the matter is already
decided beforehand: that is as unparliamentary
as possible. Even if the abuse does not go to such
% --- p. 52 ---
\setcounter{footnote}{0}%
\opage{52}an extreme, even if there is in the assembly
sense and receptiveness for grounds, there is yet ---
if true political experiences are lacking --- neither
in the moral nor the juridical factor any security
that it will become more than sham grounds. It is so
tempting for a gifted parliamentary speaker to apply
means that dazzle, to shine by wit and humor, to
wound his opponent with the poison-drop of satire and carry
the assembly away with him. Whether what he defends is in the
best agreement with his own past and will
hold in a time to come does not trouble him:
he is bewitched by the honor of the day, and relies on the public % print: Hæder. (misprint, stop for comma)
being forgetful. Yes, the public is forgetful; it
has, as Karker says of Haakon Jarl, ``so many things
in its head.'' But a people with political experiences is
not forgetful; it does not let itself be infatuated by a moment's
fireworks, but remembers precisely that on which it
depends, remembers who it was that spoke, and what
was said, and remembers it long.

To rely on the moral factor without securing
the juridical (constitutional-law) one is impractical; but to
rely on the juridical alone and to scorn the
moral is in truth also impractical. Those who have
the floor in a parliamentary assembly must know that they
have not merely listeners, but also judges, serious
judges, who in the course of time will call them to
account for every improper word they have spoken.
Only when it comes so far that a parliamentary speaker
who has won a momentary victory, which he
knows, be it ever so brilliant, rests
on hollow grounds, must say to himself what a Greek
% --- p. 53 ---
\setcounter{footnote}{0}%
\opage{53}Prince said to those who congratulated him: ``if we win one more such
victory, then it is all over with us,'' only then is
parliamentarism carried through. But then it will also
be seen that it rests on more than a juridical mechanics of voting,
that its foundation is a dearly acquired
treasure of true political experiences.

% --- p. 54 ---
\setcounter{footnote}{0}%
\lecture{IV.} {Free Constitution and Freedom.}
\opage{54}It is now 28 years, thus nearly a generation,
since a new age began for the peoples of Europe.
It was in '48 that the spirits of life stirred: the morning
shone, the day broke forth, and the sun of freedom rose over
jubilant millions. From Paris, from Vienna, from Berlin,
from Copenhagen, everywhere the same message sounded. Where
the old bonds could not, as with us, be loosened with
gentleness, there they had to yield to force, there they had to
burst. What had long been fermenting in possibility, in the
womb of disappointed expectations, of dim, unfulfilled wishes,
came at one stroke to reality; and the reality
was surprising. In one place people knew
nothing of what was going on in another, and yet everything happened
as if by secret agreement. The morning came, and behold!
the edifice of the new freedom stood complete before the gaze of all,
like Aladdin's castle, a work of the spirits of the lamp, conjured
forth in a single night.

But the enchanted castles of free constitutions have unfortunately
not been able to satisfy everyone in the long run, in particular
not those who, without troubling themselves about the style of the
% --- p. 55 ---
\setcounter{footnote}{0}%
\opage{55}whole, think only of putting their private household affairs
in order. There is in the new edifices of society
a good deal that is not pleasant. They are
magnificent enough to look at, but there is much false gilding;
they are airy enough, but there is so much
draught; they are roomy, indeed far too roomy. There
is a noise and unrest, a bustle and traffic, a shoving,
pressing and running up the stairs and down the stairs,
not only on the main stairs but on the back stairs,
indeed, what must especially annoy those who
like peace in their own kitchen, even on the
kitchen stairs. The ground is not firm, the pillars not
strong. The floor rocks; the ceiling cracks; one feels
as if in a perpetual earthquake. Every moment
something comes tumbling down; one goes in mortal
fear of getting a beam, a Rigsdag member or
a minister on one's head. One consoles oneself that the building
is so new that it will surely hold; a wretched consolation!
One has examples, after all, of new buildings falling,
even before they have been inhabited.

No, then there was a different order in the edifice of the state
under the old absolute kings. Every estate stood
at its own height; one knew how to make a distinction between
the gentry on the beletage and the servants in the mezzanine.
Now the lowest runs together with the highest:
the lady on the first floor and the woman in the cellar, the Born and
the Unborn, hat and cap, uniform and smock are all one and the same. And then people mock the Table of Ranks and
make a fuss about the Constitution. How perverse! The Table of Ranks
is the true Constitution, the only
law that can separate the classes and keep order in society:
when the Table of Ranks is abolished, the state must
% --- p. 56 ---
\setcounter{footnote}{0}%
\opage{56}go to ruin. It is important for a civilized society
that it know its great men; but how is it
possible to know them when there is no Table of Ranks?
When in earlier days one met a man with
a white ribbon and a star, one was seized with involuntary
reverence, one felt that this was a great man,
a man deserving of the state; in what his merit
consisted one did not need to know, one had after all
the Table of Ranks. Even a little cross was a gift of grace,
thus something different from a plain and simple gift of nature.
Distinctions are just as indispensable for the sense of
society as comma, semicolon, full stop and exclamation mark
are necessary for the sense of a book. But
the liberal constitutions overlook the distinction. Nowadays
deserving men go about and look like
other men; that gives occasion to confusions.
If one wants to know the deserving man, one must oneself
take the trouble to investigate in what his merit
consists, and then the yield of the inquiry may still
be thin enough. On oral traditions one dare
not rely: one must go to written sources, that
is to say, to the newspapers; but the newspapers? One paper
presents him with a halo round his brow, the
other with big blots of ink on his face. What is one to
think, what is one to believe? Might one not
wish that the political fairy tale of the people's stir would end with
a bang like a fairy tale, that Fate or Chance or
Heaven or the Lamp would send us a statesmanlike
Nureddin, who could move Aladdin's castle, this castle in the air
of freedom with its false princess, over
to Africa.

% --- p. 57 ---
\setcounter{footnote}{0}%
\opage{57}Such expressions of involuntary ill humor would
certainly have something to be said for them, if the standpoint
from which they proceed were not entirely mistaken. But
the fact is that the free constitutions on which the present
is working are not products of the arbitrariness of the multitude,
but a work of historical necessity.
The great stages of development in the history of mankind
may be compared to the classes in a large school.
When the pupils have sat a certain time in one of the lower
classes, they must either be moved up into a higher or leave
the school. The nations that, after having passed through
the intermediate stages of European civilization in the
17th and 18th centuries, are still not ripe
to take hold of the work of freedom, must put up with
going out of history. Over the stage in
our Royal Theatre there stands an inscription: \emph{Not for Pleasure Alone};
this inscription ought to stand in large golden letters
over our new edifice of society.

A playhouse, however excellent its style and appearance may be
otherwise, is not erected for its own
sake, but for the sake of the drama. Were there no
dramatic poets and no actors, the
magnificent building would stand as an empty monument, a
memorial of want and spiritual impotence. This holds in a
still higher sense of the constitution founded on freedom.
If there are no ideas of freedom and no men
who understand how to bring them forth, then the building stands
there like a temple without gods and looks scornfully down
upon a degenerate generation whose spiritual marrow is consumed.
The free constitution exists for the sake of freedom; it is
a foundation of freedom with external conditions, but not
freedom itself. A free constitution need not always
% --- p. 58 ---
\setcounter{footnote}{0}%
\opage{58}be bought dearly, it can after all be had as a gift: we
thank King Frederik VII for his gift. But freedom
itself cannot be had as a gift; it must be bought at
the dear price, it must be acquired in struggle, though in
peaceful, bloodless struggle.

That freedom must be acquired lies in the nature of the case;
for what is freedom? The first thing of freedom is self-determination,
and self-determination is activity, in the innermost
and deepest sense self-activity. The human being
who does not feel ability and strength for self-determination is
unfree, however fortunate his external circumstances may be, and the
same holds of a whole people. A people may have
a free constitution and yet find itself in a condition
of unfreedom, either because it has not yet learned to
use freedom or, what comes to about the same thing,
because it misuses it. To make rational use
of the external conditions of freedom, in particular of popular freedom,
two things are required: forces and
insight. Forces without insight give only commotion, and the
blind, thoughtless commotion makes the people giddy. Insight
without forces gives only doctrines, and empty doctrines % print: Doctriner. (misprint)
are powerless in the face of reality.

When the crises of freedom begin, the popular spirit awakens
and asks for forces. There are those who
preach and say that ``a people's spirit'' need not look
for forces, since it has the forces within itself,
and can, though we be ever so enfeebled, come
and give us all the forces we need, whatever day it
pleases; this is a bungling with spirit and words which misleads
instead of guiding. Should the spirit bring
the forces with it from heaven, they would have to be
supernatural forces, and the upswing of popular life would rest
% --- p. 59 ---
\setcounter{footnote}{0}%
\opage{59}upon a miracle. That is superstition. In a commercial crisis
the merchant says: I must have money, wherever I am
to get it from, otherwise I go bankrupt. In a crisis of freedom
the spirit says: I must have forces, wherever I am
to get them from, otherwise it is all in vain.
In this respect it proceeds just as naturally in a
crisis of freedom as in a money crisis. Yet here there is undeniably
also a difference. Capital can hide so
cunningly that the merchant is unable to hunt it
up; but forces cannot hide thus
from the spirit; it will certainly know how to find them, wherever they
may be. It is not always certain that they are found in the uppermost
stratum of society; there is a civilization that consumes
the forces, there is a sunburn of culture that
dries out the mind, so that the germs of life wither. If the
spirit finds the upper strata burnt out, it works its way
deeper down and finds the forces in the lower. The popular spirit
performs no miracles, it does not create out of nothing;
but it uses what is given: its power consists in
awakening the slumbering forces and raising them through the
growth of spiritualization.

Here is a standpoint from which the spokesmen for
formal intelligence think they can strike a blow
in honor of what they understand by enlightenment and education.
If the spirit, they say, is to seek its forces in the lowest
strata of society, then it must content itself with crude
forces; how then do matters stand with forces
and crudeness? If one assumes that only the crude forces
are unweakened, then one must above all seek to preserve
crudeness in order to preserve the forces. If popular freedom is
conditioned by forces, and the forces conditioned by crudeness,
then the age of freedom must necessarily become an
% --- p. 60 ---
\setcounter{footnote}{0}%
\opage{60}age of crudeness. If one assumes that force can be separated
from crudeness, so that the forces can be preserved after
crudeness has vanished, how then will one be able to
scorn the highest products of civilization,
art, poetry and science? If the age of popular freedom is to
be an age of enlightenment and education, then the mythological
twilight of half-enlightenment is not sufficient to
displace the raw air: the sun of intelligence must stand
shining in the clear sky. But if one admits
that the forces are not weakened but ennobled by enlightenment,
then the noblest forces must surely be sought in the uppermost,
sunlit stratum of society, and by no means lowest down in
the darkness of crudeness.

It is the well-known dispute over cultural intelligence
and popular enlightenment that we touch upon here. That art, poetry and
science contribute mightily to awakening and clarifying the life of the spirit,
and that the forces of freedom, so far, must rather be
sought in the upper strata than in the lower, is a fact
which no one who knows a little of history will be able
to deny. But that culture, besides this, hatches
a swarm of infusoria of intelligence, who reason
on the surface and consume the forces, no one else will be able
to deny either. Superficial reasoning goes from doubt
to doubt, and the endless doubt weakens the power to act.
What good does it do a civilized present to know all
the virtues and vices of the past, if it does not know
and understand itself; and what use is it if it
also contemplates itself in the mirror of intelligence, so long as it
remains with the barren contemplation. Light without forces
is a dead glow; but where the forces are, the spirit can
find them and kindle light, even in the deep darkness.
% --- p. 61 ---
\setcounter{footnote}{0}%
\opage{61}The word of truth does not say that the light is the life of man,
but it says that the life is the light of man.

Where strong forces give birth to clear thoughts, there comes about
an education by which what is common to humanity
and what is peculiar to each capable person
is developed. In the popular society the development of freedom
will then show a multiplicity of stages of education and circles of education,
and the education in every circle will correspond exactly
to the labors that are assigned to the particular classes.
It is quite true that the hardy men who move about
on land and sea understand better how to handle
reality than those who sit in a study
and busy themselves with thoughts; but for that reason it ought not
to be forgotten that there is also a labor of thoughts, that
many a light which has shone over society and set
a world of new forces in motion has been kindled in
the quiet chamber. The main thing is that the different
kinds of education enter into the right reciprocal action,
in that they all, each with its measure of forces and light,
enter the service of the development of life and of freedom. The core
of all education is self-education, formation of mind, the education
that teaches the individual that he in truth only
works for himself by working for the whole. On
these conditions the popular spirit can stir, and the work of freedom
be advanced.

The forces can work without spirit; but the spirit
cannot work without forces. Where the forces work
without spirit, society will find itself now in drowsy
order, now in uncanny ferment. Either the activity will
be so distributed that each estate attends to
its own, the government to its own and the people to its own; here reciprocal action is lacking.
Or the opposite endeavors will fall
% --- p. 62 ---
\setcounter{footnote}{0}%
\opage{62}into strife with one another, one party will combat the
other and lose sight of the common aim: here there is lacking
harmonious cooperation. The mark that it is
the popular spirit which sets the forces in motion is this
twofold: that all endeavors work together toward one
aim, and that all, high and low, are conscious that
they must take part in the activity. The latter point
becomes especially important when the free constitution has
succeeded absolutism.

In a condition of society where the spirit is still
half slumbering, it is natural that only some
individuals of particular classes concern themselves with public
affairs, while the great multitude keeps to
the private and lets the government look after the state. It
is a view which pleases the government; it does
not care for restless heads. The authorities seek, as
guardian of the people, to establish firmly the principle that
the subjects have to submit to their measures
and for the rest keep quiet.

It lies in the nature of the state to maintain itself as
something existing; those who govern therefore keep a sharp
eye on the restless, but look graciously and mildly
upon the quiet in the land. The statesmen of antiquity were
in this respect already attentive to the difference
between the workers in the cities and the rural population.
``The worker,'' says Aristotle, ``is a man
of small virtue; and only in imperfect states can it be
permitted to workers to take part in legislation and government.
A population that is chiefly occupied
with agriculture and cattle-breeding is as a rule content,
if it can merely get peace to continue at its work,
and is willing to content itself with a limited share in
% --- p. 63 ---
\setcounter{footnote}{0}%
\opage{63}the administration of the state. Far more restless is the mass of craftsmen,
tradesmen and such as live by profit:
their way of earning a living acts more harmfully on
character, and in the city, where they are found crowded together,
they are inclined to make noise in the assemblies.''

When it is to become serious with the use of the
free constitution, that is, with freedom, we cannot --- with all
respect for peace and order --- praise the
quiet in the land unconditionally. What kind of people are these? Decent people;
in private, personal relations honorable, respectable,
useful members of society. But their inclination to
peace and zeal for minding their own is so predominant
that active participation in public life is almost
repugnant to them. It is in vain to remind them
of the inscription: Not for Pleasure Alone, for they do not even
have pleasure. Nor is it of much use to set before
them the common duty which the constitution imposes on us in
the name of freedom, for as to freedom, they have
their own explanation: ``In a free constitution everyone,
who does not wish to join in the public bustle, must surely
have freedom to hold back. There are, they say,
people enough who have a mind to act on the political
stage. But there is no real pleasure in
being an actor when there are no spectators. We, who
cannot act, serve the public by being
spectators.'' The quiet in the land are spectators. What
joy do these spectators now have of the popular-political
scenes? The joy is not great. Those who sit in the first
rows of the pit and can plainly see everything through the spyglass of criticism,
do nothing but fret. If the play goes
briskly, they fret, for freedom becomes saucy and
throws sand in their eyes. If the play drags, they
% --- p. 64 ---
\setcounter{footnote}{0}%
\opage{64}fret also: there ought surely to be hissing; but hush! ---
for the sake of quiet one scarcely dares even to sneeze,
still less to hiss. Those who stand so far back
that they can scarcely make out what is going on on
the stage regard the whole as something that basically does not
concern them; the only thing that can occasionally give
their tranquility a jolt is rumors of some frightful
things that are said to be going on behind the scenes.

Since the quiet in the land do not express themselves publicly,
we cannot of course have any reliable
knowledge of their opinions in detail; but that
may also be all the same. What matters here
is solely the fact that they do not regard
an active participation in the great affairs of society
as a general civic duty. Could we directly address
questions to them, they would doubtless answer:
``Nothing comes out of all this political rummaging.
It cannot be every man's duty to understand
affairs of government. We do not understand
affairs of government and consider ourselves too good to do
drudgery for those who, in order to be able to use us,
must begin by making us imagine that we are great
politicians. A sensible man never concerns himself with
what he does not understand.'' --- Such an answer would be
adequate under absolutism, but under a free
constitution it is by no means adequate. If
a human being could without more ado get out of concerning
himself with what he does not understand, then it would surely be
clever to apply oneself to understanding as little as
possible. There may be much that we do not understand, and which
we yet must learn to understand, when it concerns us
and our work. Certainly not every other person in
% --- p. 65 ---
\setcounter{footnote}{0}%
\opage{65}a population can sacrifice time and diligence to know
the affairs of government in detail; but if we put the
question to the reserved, whether they then would give up
their franchise? we shall doubtless get the answer: ``by no
means! why should we not have the same right as
all the others.'' --- But how can you set a price
on a right which you, by your own admission,
do not know how to use? --- ``We understand it just as well
as all the others.'' --- Then tell us, how you
go about it on an election day? --- ``Sometimes we turn up,
sometimes we do not turn up. But every time we turn up,
we get right into the matter. We hear the candidates'
instructive statements and take note of the enlightening answers,
when they are interpellated. To be sure, we cannot always
find our way so readily in the details, but the spirit is after all
the most important thing.'' --- How do you go about
grasping the spirit? --- ``That is the easiest thing
in the world. One need not, for that matter, even
understand what it is that is said, but merely take
note of how what is said affects the assembly.
If it elicits Hear, Hear! and Yes, Yes! then the spirit is there;
if it elicits Hush, Hush! No, No! then the spirit is not there.
If we now find, after careful consideration, that the spirit
is for a man of the Right, then we naturally elect a man of the Right;
but if the spirit absolutely wants a man of the Left,
very well! then we can also elect a man of the Left.''

--- So: in this manner you serve freedom by
conscientiously! --- using your franchise.

The franchise is a dangerous weapon; to entrust a dangerous
weapon to people who cannot even be bothered to inform
themselves of its right use is thoughtless. There is complaint
that we have too many examinations; one could after all
% --- p. 66 ---
\setcounter{footnote}{0}%
\opage{66}be tempted to propose one more examination, namely a
general voters' examination. How many commissions
would have to be appointed on this occasion, how
they should be composed, of which and how many
members such an examination commission ought preferably to
consist, how large per diems, in proportion to the price of grain
or the tobacco tax or the rising export
of eggs, might accrue to each member of the commission,
these are practical considerations that we cannot here go into;
it is enough for us to point to the need of the time, its
deeply felt need of a voters' examination. In such an
examination it ought not chiefly to depend on
the candidate's external position, whether he is common or
distinguished, poor or rich --- respect of persons does not count
in a public examination. --- The outcome depends on
the candidate's answers, and on order being kept with the marks.
If a candidate can give sufficient account of
his point of view and of the reasons which guide him when he
elects, then it ought not to be to his detriment that he
is for the rest a man without means; even a small
running account with the poor relief shall not prevent
him from getting a Very Good with a cross. If a candidate,
on the other hand, presents himself without proper preparation, weaves
in and out about Conservatives and Liberals, about the party of standstill
and the party of progress until he
literally cannot tell Right from Left: then the
commission ought not to let itself be dazzled by his external advantageous
position; even if he owned ten factories and
twelve manors, he ought nevertheless to have a mediocre
mark: for my part I would vote for Zero.

But --- one will perhaps object --- such a proposal
will never pass; it runs into practical
% --- p. 67 ---
\setcounter{footnote}{0}%
\opage{67}difficulties. Possibly so. But for that reason we shall
not lose heart. We have in reality, when
we look more closely, a voters' examination. Existence itself
with its experiences is the great examiner. Freedom
has its rod of discipline in strict necessity; it
disciplines us by the consequences of our actions. A free
man, of course, is permitted to sleep. But in the
busy season, when the seed is to go into the ground or when the harvest
is on, the free man knows very well what consequences
it may have to oversleep. As regards agriculture and
industry, a free man is certainly permitted
to stand still and say: I do not keep up with the times; with
me everything shall stay as it was. I act as
my grandfather acted, when there were as yet neither
steamships nor railways. I keep my farm in the
same operation as my forefathers; the newfangled inventions
do not concern me; I will not depart from my
old habits. See, that a free man can certainly
do; for the free constitution will hardly issue any compulsory law for the redemption of lazy habits; but the consequences
do not fail to appear: freedom will discipline the free man
with the rod of necessity. Thus, and not otherwise,
will it go in the course of time with the use of the franchise;
one can misuse it and one can neglect
to use it, but the consequences will not fail to appear.

Those who complain that things are going awry with freedom,
and by their complaints imply that they know better
how it ought to go, dare not sit by as spectators.
They must come forward on the stage, that the public
may see the coat of arms they bear. Freedom is activity,
self-activity, and self-activity is closely
akin to enterprise. The enterprising party in a country
% --- p. 68 ---
\setcounter{footnote}{0}%
\opage{68}will always be a horse's head ahead of the merely reasoning one.
The encroachments of an opposing party cannot be repelled
by merely pointing out its shadows and crookednesses, but
by organizing, closing ranks and acting.
To be sure, freedom may also demand of its resolute ones
that they for a time stand still, see how things go and let
the hot-tempered opponents run out the line. But such
a standing still is then itself an action, and such an
action widely different from being either a sulky
critic or a thoughtless spectator.

To invoke the popular spirit while yet doing
everything to stifle it is certainly a contradiction;
it is a sad sight when an enterprising party of progress
in blinded defiance chooses the wrong means. So long,
however, as there is faith in the spirit and courage to
fight for the cause of freedom, there is yet hope; the mistakes
can surely be corrected, and the development again come into the
original track. But those who are in the process of
giving up the spirit, the circumspect men of order, who
are calm only because they are weary, and weary, indeed dead tired,
because their capital of life is used up, because they
have not a single viable and fruitful thought
left, what hope is there for them? Their watchword
is: to keep \emph{things going}. Of course, to
keep things going can also be a task for those
who are in full activity with what is future. But
when keeping things going is the last refuge of weariness, of the burnt-out
and yet of tenacity of life, then
there falls over it a glow of old age which is uncanny.
There is told in Snorri Sturluson of a King
An or Øn, how he sought to keep things going.
It says in Grundtvig's translation: ``After Jørrun
% --- p. 69 ---
\setcounter{footnote}{0}%
\opage{69}his son An or Øn became king in Sweden; he was
a clever man and a great idolater, but war he
left to others, and held home dearest. To his
misfortune it was otherwise with the Danish
King Halvdan\dots{} He went to Sweden with an army, and
fought victoriously with King Øn, until
he had to take flight into Vester-Gylland, and let
Halvdan rule till his dying day in Upsal. At that time
Øn had already reigned twenty years, and Halvdan reigned % print: eg (misprint)
another twenty, but when he was laid in the mound, Øn came
back to Upsal in his sixtieth year, and made
at once a bloody sacrifice, in order to live long yet. It
was his own son whom he sacrificed to Odin, and received for it
the promise that he should become just as old again as
he was, but the kingdom he did not keep this
time either longer than twenty years, for then came King Fridleif's
son in Denmark, Ole Bravkarl, with his army, struck
Øn both asunder and together and forced him again
to flee to Vester-Gylland. Now twenty years passed again,
but then Stærk-Odder killed that Old Ole Bravkarl,
and at once Øn came back to Upsal, and ruled
his twenty years for the third time. He wanted, however, to live
longer yet, and sacrificed to Odin his second son, for which
he received the promise to live ten years, and so on ten
for each son he gave Odin, but on the condition that
each time he should give a district of the same number
as the son he sacrificed. Now he went on, and when
he sacrificed the seventh son, he could no longer
stand on his legs, so that he had to be carried in a chair,
but nonetheless he sacrificed the eighth too and
could in the following ten years not get up from his
% --- p. 70 ---
\setcounter{footnote}{0}%
\opage{70}bed. Still he sacrificed his ninth son, and had,
to keep his life, to suck like an infant at a horn.
One only son he had still, and him he would
also have sacrificed to Odin, together with Upsal and the district
around it, which should have been called Tenth-Land, but this
the Swedes denied him, so that nothing came of it, and
An had to resign himself to die, and be laid in a mound
at Upsal''.

As there can be a wondrous tenacity of life in
old worn-out people, so there can also be an
almost incredible tenacity of life in old, worn-out states of affairs.


But not merely laxity, will-lessness and the debility of old age
is a burden upon freedom; the
unruly enterprise as well, the selfish pushiness, can
put it in danger. --- A foreign potentate is said once
to have said: every time Denmark is about to go
under, the Lord always has a little miracle in reserve.
We have no right on this point to rely on miracles,
but we are bound to rely on the aid of the spirit,
heed the signs of the times and acknowledge the claim of freedom.
Fresh forces, fruitful enlightenment,
high-minded courage and popular education are ideal powers,
with which we must first and last seek alliance. They
are allies who will not fail in the hour of danger, if
we only do not fail ourselves. But we fail ourselves
when we pettily squabble about power; we give ourselves up
when in passionate blindness we transform our constitutional edifice
into a gambling house, where the party spirit says
va banque and stakes freedom.

Still stands Denmark's old house, despite robbery and
plunder, on its own ground. Of the timbers of freedom
% --- p. 71 ---
\setcounter{footnote}{0}%
\opage{71}shall the small remnant that is left be built up anew.
It will presumably not become an Aladdin's castle, but if the workmen are
ready, and the master builders agreed about style and plan,
the house can rise with the wreath of honor, a memorial of concord,
a work of the spirit, a stronghold of the people against the storms
of the age.

% --- p. 72 ---
\setcounter{footnote}{0}%
\lecture{V.} {Aristocratism and Democratism.}
\opage{72}The busy present must never forget how much it
owes to its time of preparation, to its honorably buried
past. Should it forget this elsewhere, in this
town it will not forget it: the hand of the past is stretched out here
and reaches into the present. Youthfully fresh with
hope and outlook for the future, life can stir everywhere,
in the houses of the citizens, in the workshops of the workmen, in
the squares and streets, but over the bustling life of the present there
spreads quietly a gleam of antiquity, an illuminating light from the
vanished past: Viborg stands on historic ground with
old thoughts and with old memories. The cathedral yonder
with its round arches and its high towers keeps the
memories; the school opposite, in its sight, gathers up the
thoughts, and both, church and school, agree that
everyone who wants thoroughly to understand the intricate present
ought not to neglect to live himself into the great past, so rich in experience.

Past and present, the old and the new, are
in unceasing conflict; a ceaseless striving after
change is forever at work transforming what exists.
The old, it is said, must fall and the
% --- p. 73 ---
\setcounter{footnote}{0}%
\opage{73}new rise on its ruins, only progress is life,
standstill is death. In such political catchwords there is
something true, but also something untrue, awry and misleading.
The question is how things really stand in society
with the transition from the old
to the new.

When a house grows old and falls into decay, one repairs
it. One can then build up much, take
away something and add something; but when this happens
again and again, the building too can gradually become
so altered that the original style can no longer be
recognized: the partial transformation ends in a baroque
patchwork. The English structure of the state, with the
reforms it has undergone in the course of time, has much
that reminds one of this. An English parliamentary speaker
once remarked, on the occasion of the multitude of
alterations that were being made in a Reform Bill:
this Bill reminds me of the cotton stocking that was
darned so long with silk until it became a silk stocking.


The opposite procedure is applied,
for example, in France, where one briskly
pulls the old down to make room for the new.
When what exists no longer satisfies, one makes
a revolution. In the course of a couple of days one can,
when one has got properly into practice, pull down a whole structure of the state
and introduce a brand-new constitution.

In the first case --- after the English pattern --- it is
the old that partly distorts the new; in the
latter case --- after the French pattern --- it is the
new that destroys the old. But precisely through the
% --- p. 74 ---
\setcounter{footnote}{0}%
\opage{74}old being destroyed, the foundation of the
new is undermined. The latter is worse than the former.

The relation between the old and the new
takes a wholly different shape according as it is a merely mechanical
or an organic, a lifeless or a living building
that one has to do with.

The human structure of society is not a
dead mechanical composition, but a living, animate
organism.

In living organisms, plants and animals, there goes on
an uninterrupted transformation of all, even the
smallest parts. Work is done in cells and tissues and countless
workshops on the transformation, so that the change
becomes thoroughgoing. The basic form, what
we might call the style, is maintained. In periods the
basic form too may indeed be altered, --- there are, as is well known,
plants and animals that in different periods of life
have an entirely different appearance --- but even where
the change seems to happen suddenly, it has
long been prepared, on all sides and in an inward manner, so
that in all members and parts there appears the most intimate
connection between the old and the new.

What nature has set before us in examples from the
living organisms should now repeat itself, in a higher degree and in a
higher manner, in human society.
But precisely because the conditions of development are
of a higher kind, the development itself undergoes various
oscillations.

Social life with its many institutions does not
shoot forth like the shoots of a germinating seed, does not grow in a
quiet growth according to necessary unalterable laws; for
development has a twofold basis, a natural and
% --- p. 75 ---
\setcounter{footnote}{0}%
\opage{75}a spiritual one. The laws of society, however
much they may be necessary by nature, are laws of freedom, with a stamp of
arbitrariness; development is a struggle between many
interests. If the struggle becomes so violent that the strata of society
topple over one another, that class stands up
against class, party against party, then the even development
is interrupted, and the question becomes how the
new will place itself in relation to the old. No
people leaps with impunity in its historical development. Traditions,
habits, all the after-effects of the past continue
with a terminal velocity that is decisive
not only for the speed of the movement but, whithersoever
the new may otherwise strive, also for its direction.
The transition from one constitution to another is
always precarious, and that in the same degree as the corresponding
conditions are unlike. A people that for centuries
has lived under a monarchical government is ill
served by being flung headlong into a democratic
republic. After the experiences that peoples and statesmen
must have had in the latter half of the last and
the beginning of this century, one would suppose
that the time must now be past when people imagined
that ready-made theories of the state could be obtained in
a political shop, like ready-sewn clothes in a
tailor's warehouse. None the less the party spirit,
when it finds its account in it, without regard to
the organic conditions of life of the people concerned, will naturally
not fail to beat the drum in urgent
recommendation of democracy. Democracy is a
catchword that always suits the taste of the multitude, and why
should the demagogues not have it in their mouths?
When it is said of a man that he is an aristocrat,
% --- p. 76 ---
\setcounter{footnote}{0}%
\opage{76}then the mob knows at once that he is an unpopular man,
a man who overlooks and despises the people. Since the
aristocratic in this way becomes synonymous with
the unpopular, the democratic naturally becomes
synonymous with the popular. The democrat of the pure
water, the genuine democrat, is the friend of the people; so the multitude
believes, and takes the catchwords in use on this occasion
for good wares. Natural enough, when it does not
know the mixture of true and false that lies in
such catchwords, when it is, in other words, not
able to take counsel with history and to feel
its great lesson. It is undeniably true that there
is a form of aristocracy --- be it aristocracy of birth
or aristocracy of office --- that is incompatible
with the demands of popular life. But the same, the
very same, holds of democracy. In democracy too,
in radical democracy especially, there is something
unpopular, something worn out and threadbare, that must be transformed.
We can best convince ourselves of this by
looking at the principles.

What sort of principle is it, then, on which
aristocracy stands firm, when we look away from the historical
abuses? it is literally and according to the wording this,
that \emph{the best shall rule}. If the democratic
were now what it gives itself out to be, the direct opposite
of the aristocratic, then the principle of democracy
would evidently have to be this, that \emph{the worst
shall rule}. But such a principle
surely no one will acknowledge. If the democrat will declare
that it is neither the best nor yet the worst
who shall rule, but the multitude, and that in the multitude
all are equally good: then we must ask him what
% --- p. 77 ---
\setcounter{footnote}{0}%
\opage{77}it is then the multitude looks to when it chooses a
leader. If it is really true that all are equally
good, then it is surely not necessary to make particular
fuss over the election, one can choose blindly or
by drawing lots or however one likes. But although
the multitude knows, or at any rate imagines it knows,
what ought to be done, yet when it comes to the point,
it is not everyone in the multitude who knows how
it ought to be done, indeed, if one looks more closely, it is
only very few who have any tolerably clear idea
of what it really is that ought to be done. The multitude
has by instinct a feeling that although all
are in general equally good, yet
all, when it truly concerns carrying out something,
are not equally good. It is of importance to the multitude to get
capable leaders, for without leaders it can accomplish nothing.
When it has found a leader who is entirely
after its mind, what does it do then? then it follows
him blindly; it obeys him unconditionally, believes in
him --- that is to say so long as it lasts --- and why?
naturally because it regards him as the most capable,
the best. The best is therefore to lead, the
best is to have the power, the best is to rule.
But that the best shall rule is surely an aristocratic
principle. The multitude is therefore, without knowing it,
aristocratically disposed.

It is incalculable what popular life would gain,
if this simple principle, evident to all,
that the best shall rule, were allowed to make its way
through.

Let us consider the opposition from another side
and look not merely at the persons but at the cause, not
% --- p. 78 ---
\setcounter{footnote}{0}%
\opage{78}only at the best, but at the best thing, the true,
the good, the right, in short: that which serves the whole of society.
How do matters now stand between
the aristocratic and the democratic? The aristocrat
presupposes that civil society, whatever
the constitution may be, is subject to definite laws of development,
that there is something harmful in and of itself and something
beneficial in and of itself, and that the difference between harmful
and beneficial cannot be decided by loose judgment or
by a momentary majority vote, but only by investigation of the matter
and insight. His motto is that the truest
is always the best. If he imagined
that one need only be born with a manor
or an entail in order to have a privilege
of knowing what is best for society, then
his aristocratism would certainly be a delusion; but
the true aristocrat has a respect for truth and
right that he cannot give up without giving up himself.
--- --- Let us then see how the democrat goes about it.
To find out what is for the best of society, he
need not trouble himself with laborious investigation;
he has only to heed the opinion of the most;
``the opinion of all,'' and to turn to the multitude.
The multitude is his oracle, the multitude alone knows what
is the best. No pointer has a finer nose for
tracking down game than the democrat for tracking down
the opinion of the multitude. The acuteness he thereby
displays has a fair claim on our admiration.
In the first place he has an opinion himself; but to acknowledge
it and step forth in public with a wretched
individual and private opinion, what could that lead to?
It would surely be subjected to sharp criticism; for
% --- p. 79 ---
\setcounter{footnote}{0}%
\opage{79}one pays no compliments to an individual opinion;
indeed the public might even be capable of demanding
of him the impossible, that he should in earnest
engage in giving grounds for his opinion. No thank you!
He is not so stupid as that. He chooses a more practical
procedure. As he steps forth before the multitude,
he knows how to speak in such a way that, by an acoustic
deception, --- for the democrat is a practiced ventriloquist ---
it comes to sound as if it were the multitude's own
opinion that he was uttering. The multitude, which now for the first time properly
learns that it has an opinion, at once proclaims
the opinion far and wide, and, as gladly surprised
at hearing its own opinion as the woman in ``Gjenboerne''
at seeing her own ``winders,'' it shouts to the skies
in rapture. Now there remains only, what is
the easiest thing in the world, to lump the multitude and
the people together, and what a transformation! The democrat's
opinion has become the multitude's opinion, the multitude's
opinion the people's opinion, and the people's opinion needs
no justification.

On these presuppositions the democrat can defend
himself before any tribunal whatever. It is
in vain that the experts come forward with objections;
he answers merely: Thus the people want it,
and one must bow to the will of the people. It is in vain
that the better informed propose something that
in itself is far more reasonable; if the democrat cannot
defend himself against the telling reasons that speak
for his opponent's proposal, then he says: As a doctrine
the matter may well look all right, especially when it
is seen through learned spectacles; but the people, who prefer to see
things with their own natural eyes, have once for all
% --- p. 80 ---
\setcounter{footnote}{0}%
\opage{80}come to distrust this matter, and to offer the people anything
it distrusts is impractical. One sees here
the fineness of the procedure: first he awakens in the % print: Fremgangmaaden (misprint)
little expert people a distrust of the matter he does not
want brought forward, and then he applies the distrust thus awakened
to get the matter rejected, before the people have yet
come to see their error. Thus must
the genuine democrat go about it when he is properly to
act in accordance with the democratic
principle.

But, it will be asked, are all democrats then really
such deceivers? We answer no, far from it, the most
among them are without doubt honest, well-meaning persons,
who stand on the ground of freedom and really
imagine they are serving the popular cause. The secret
is that they read something else and something more into the principle
than it actually contains. The principle
is that the multitude knows what is the best, and
that the will of the majority should for that reason be the only one
ruling. The deception hidden in this comes to
light only when one gives a serious account of \emph{truth}
and majority.

To arrive at knowledge of truth --- this holds
of political truth too --- is, and has from time
immemorial been, a very difficult matter. Not the multitude, but
only quite few were granted to see such clear glimpses
of the light of truth that they found a way through the wilderness % print: igjennnem (misprint)
of error; still fewer were granted to press
through and win an entrance for truth. In antiquity
truth was regarded as a secret that might be entrusted
only to the initiated; to communicate truth to the great mob.
was to cast pearls before swine, an enterprise dangerous to life,
% --- p. 81 ---
\setcounter{footnote}{0}%
\opage{81}a violation of the holy. That was a false aristocratism.
It is true that the discovery of truth is reserved especially for
the gifted, but the truth once discovered will
not therefore be shut up within Egyptian temple walls:
truth is light and wants to shine for all. It goes with
the truth of life as with beauty; there is a natural
sense for the beautiful, and there is also a natural
and original sense for the true. The human
eye must gradually accustom itself to the glimmer of light.
A sudden transition from the gloom of night to the strong
light of the midday sun is painful, the rays dazzle and
it cuts into the eyes; but little by little the transition
may take place, and then the light is beneficial. Human beings
have by nature an eye for truth just as they have
an eye for light, a need of truth just as they have
a need of light, yet with the difference that always
obtains between the merely natural and the spiritual.
The spiritual eye and the spiritual need must be
awakened, purified and sharpened, that is to say, there must be
education, if the human being is truly to become
a human being.

From this standpoint the opposition between the
aristocratic and the democratic is striking. Both,
the aristocrat as well as the democrat, influence the multitude,
but in a different spirit.

The democrat proceeds from the presupposition that
truth is something innate. Just as every living
human being knows what life is, and what it is to live,
so too everyone in society knows what is
beneficial and harmful to society. Truth consists
in distinguishing between the beneficial and the harmful.
But what then is the beneficial? It is surely
% --- p. 82 ---
\setcounter{footnote}{0}%
\opage{82}obvious. Beneficial is that which heightens our well-being, which
procures for us the greatest liberties and the most comforts.
To understand this there is need of no laborious
instruction. The beneficial is what we all
desire; that the desire is present is seen from the fact that the impulse,
the inclination, indeed the passion is present. In impulse,
inclination and passion the will of the multitude resides.
The multitude knows well enough what it wants, when it is
told. On this presupposition the democrat
too knows what he wants: he wants naturally to let
the multitude know what it is that it itself wants, the
true, the serviceable, the beneficial, the best. His
speech does not aim at convincing, but at
persuading. To convince is laborious and, when one
speaks to the multitude, exceedingly tedious. But to
persuade is both alluring and enlivening. To convince
there is need of reasons; but reasons often have,
like the doctor's pills, a bitter taste, and so they lie
heavy in consciousness and are hard to digest.
To persuade there is no need of real reasons, but of
catchwords; catchwords act in the assembly like kettledrums in
an orchestra, enlivening to the mind, musical to
the eardrum.

The aristocrat too must address himself to the multitude,
for in the multitude the people is outwardly present. If he
holds himself back in silence and imagines that the people
is nothing but multitude, that one may despise
the people because in a certain sense one can overlook
the multitude, then he is only a spurious aristocrat, a
ghost from the past, with whom we have nothing
to do. Thus --- for the people's sake --- the
aristocrat too must speak to the multitude. But his speech is
% --- p. 83 ---
\setcounter{footnote}{0}%
\opage{83}different from the democrat's. The aristocrat wants more than to
persuade, he wants to convince. He pays no heed to
what pleases or displeases the multitude, but only to
that which, by his well-grounded assumption, is
the true, the serviceable and the best. He wishes naturally
that many, if possible all, would unite
around that which he regards as the best, but it
never occurs to him to regard anything as the best
because it has won the approval of many. He is not
so conceited as to regard his own view as infallible,
he is on the contrary willing to yield to counter-reasons;
but it is no use meeting him with the objection
that the majority is of another opinion than he, for
he asks not about numbers, but about reasons. The aristocratic
is to be able to stand in the midst of the multitude alone with
the truth he has recognized, bowing only to reasons,
but not to a number with many noughts.

What the majority accepts is true and valid; that
is a democratic principle. What is true
and valid the majority ought to accept; that is an aristocratic
principle. The democrat sets the majority above truth;
the aristocrat sets truth above the majority.
If it holds that a person stands the lower, the
lower he himself sets truth and right, and the higher,
the higher he sets truth and right: then it follows
from this that the democrat must stand lower in proportion
as he is true to his principle, and the aristocrat the
higher, the more faithfully he carries through his principle.
But from this one also understands what might otherwise
seem strange, that the democrat as a rule is better
than his principle, whereas the aristocrat seldom or
never quite succeeds in doing justice to his principle.

% --- p. 84 ---
\setcounter{footnote}{0}%
\opage{84}The progress of popular life will therefore depend on this, that
aristocratism steadily grows, and that democratism
steadily diminishes. The confusions in civil
society have, when we look more closely, always their origin
from above; from the distortion and abuse of
aristocratism, and, what is more, from want of ability
and will to get the abuses that had crept in abolished
before it is too late, the corruption springs. The
French Revolution affords instructive examples.

It was not for the sake of the people, but for the sake of the state,
that absolutism was raised to its height. Although the
common folk suffered unspeakably under it, yet they put up
with it, so long as the power and greatness of the state gave their
self-esteem nourishment. But when the whole power of the state had been merged
in the royal power, and the absolute royal power had
done its part to destroy the state, it followed of itself
that the foundation of the royal power had to be undermined,
and that the King personally had to bear the blame
for the public misfortunes. Leaning on the sovereignty of the King there stood
the sovereignty of the nobility with its privileges, its dishonorable
arrogance and ever-increasing presumptions; these became
the more intolerable, the deeper the state sank. The King was
the aristocrats' aristocrat, so that one cannot find it
strange that public opinion regarded ``King
and hereditary nobility as one of God's worst scourges.'' Under
the absolute kingship the Church and its hierarchy were
at bottom transformed into an institution of the state, for the King
was in reality the supreme hierarch: the royal
conscience was the subjects' standard conscience.
If now the state was ``without respected laws, its arms without
luster, the court without dignity and morals without decency,''
if besides the immorality of the higher clergy
% --- p. 85 ---
\setcounter{footnote}{0}%
\opage{85}--- one thinks of a Cardinal Dubois, a Cardinal
Rohan and others --- was on average just as unbridled
as that of the nobility, its unbelief just as conspicuous as
that of the materialists and its official profession of faith a
mask, under which it held fast only to the great
ecclesiastical revenues: what wonder if sound understanding
regarded religion as a dark superstition and its
priests as deceivers who ought to be hanged. Thus the
corruption had proceeded from above.

The danger drew near, and the blindness grew with
the danger. Already Madame Pompadour, Louis XV's all-powerful
mistress, is said to have said: \emph{after us
the Deluge}; and the Deluge came, and the waters
rose above the high mountains. There was a lack of
will, because there was a lack of understanding, a confused
conflict between old and new made the
evil worse. ``To take the shock off all this
France had only Louis XVI, a truly good but little
capable man, who, from distrust of himself, usually
sought help from still less able persons, who
were above all still less honest. He was forced
to change ministers and with them the system, and when the
bad harmed him, the good did not benefit him.
The short-sighted court, which had followed upon the boundless
corruption of the courtiers of Louis XV, did not know
how to put the King at the helm and therefore demanded
that he should stop it, but without being able to
instill in him the necessary strength. Therefore the government displayed
that mixture of arbitrariness and weakness, % print: Vilkaarligked (misprint)
which excites without cowing the desire for resistance, but
on the contrary makes this popular and gives it the prospect
% --- p. 86 ---
\setcounter{footnote}{0}%
\opage{86}of success. Tossed about among ministers, courtiers,
women, historical traditions and philosophy, Louis
XVI let the ship of state run on a half wind, and
he captivates interest only when he ceases
to act and begins to suffer'' (Cantù).

As to nobility and clergy, the mood was
divided between obstinate clinging to the old
and frivolous flirting with the new. It is so tempting
to let oneself be carried away by new ideas, so long as they do not
interfere disturbingly with one's own inborn rights
and ingrained habits. People amused themselves with Voltaire,
let themselves be inspired by Rousseau, laughed at ``Figaro's
Marriage'' and studied the Encyclopedists without
considering how dangerous it is to play with fire and light
beside an open powder tower. ``Lafayette preached a crusade
to America in the name of freedom and went there himself
to shed his noble blood for freedom. Young men, % print: Brod (misprint)
the future pillars of the French aristocracy, hastened
to fight for the overthrow of the special privileges
that remained unchanged in their fatherland, and
to learn principles of equality, of hatred of all
sovereignty of kings, ministers and priests.''

The Parliament fought with the court, and the authorities
betrayed one another mutually. The middle class had in
part joined in; but no one, with the exception of
some philanthropists, troubled himself about the common man.
The literary champions of freedom thought little or
nothing of the moral regeneration of the people; their doctrines
were calculated for the educated class and
did not concern what one then called the dregs (la
canaille).

% --- p. 87 ---
\setcounter{footnote}{0}%
\opage{87}None the less it was precisely the lower population,
the rabble despised by the men of freedom, which
through the new doctrines was to come to power. The
new doctrines had their strength chiefly in the
merely demolishing; criticism, doubt, denial were
the most effective weapons; the positive and new confined
itself essentially to bitter theories. Rash paradoxes,
boldly flung-out catchwords, flashes of thought that in a
moment lit up the darkness, flew across the masses and
set the great mob in passionate commotion. Under these
presuppositions the movement had necessarily to become
democratic.

Radical democratism is recognizable by the
elemental natural force with which it attacks and tears down
existing barriers, and by the total lack
of ability to organize anything new by developing
the germs that are contained in the old. It is not able
to separate the kernel from the shell. Whatever is afflicted
with abuse, were it in itself the noblest and
best, shall be abolished. Royal power was abused, and
so one abolished royal power; the power of the nobility was
abused, and so one abolished the nobility; the power of the Church
was abused, and so one abolished religion; Intelligence
was abused, and so one slandered science
and executed a Lavoisier. --- Terrorism is
radical democratism; Robespierre is the most radical
and most consistent of all democrats, his virtue
guillotines all abuses.

The aristocracy of birth has in its time abused its
position, the aristocracy of office has abused its position
and the aristocracy of instruction has neglected to do
% --- p. 88 ---
\setcounter{footnote}{0}%
\opage{88}its duty. By unrolling the old register of the sins
of aristocratism, history gives us sufficient instruction
about the rise of democratism. It is not political
catchwords, but a sober conception of the historical
development of society that can teach us how
the old is related to the new, what under
transformation must be that which persists, what must be
reborn and what must die away and vanish.

% --- p. 89 ---
\setcounter{footnote}{0}%
\lecture{VI.} {Politics and Conscience: Machiavellianism.}
\opage{89}Politics and Conscience! They are two great powers in our
existence. Could these two powers agree to enter into
``a holy alliance'' and keep it inviolate, what
would not thereby be gained for the good cause of humanity?
And what would not be spared? Means
and capital, painful sacrifices, costly energies and blood
and tears! But alas, the two powers still stand
so far apart that it almost seems as if they belonged
to two realms unconcerned with one another. Politics dares % print: no stop after Riger
not ask Conscience for counsel, and Conscience
can seldom or never approve the plans of Politics.
This must not be understood, however, as if the wise men
who sit in the councils of princes, the busy and clever men who
work in the cabinets, were unconscionable people
more than the rest of us mortals. Far from it! The unconscionable
does not proceed immediately from the
evil disposition of individuals. It originates in a distortion % print: Forvanksning (misprint)
of the essence of politics. Unconscionable
politics is a satanic web in which humanity---
by the guilt of thousands, to be sure---has in the course of the ages
more and more entangled itself.

% --- p. 90 ---
\setcounter{footnote}{0}%
\opage{90}It is of little use to moralize to statesmen
about disarming, about being open and honest, about always willing
what is right, and so on. They smile and shrug their shoulders,
and their eloquent silence is easy to interpret. ``Disarm,
when neighbors and adversaries are arming to the utmost?''
``Be open and honest, when friends and enemies employ
all the circumlocutions of cleverness and cunning?'' ``Will
what is right, what is in and for itself true and good, when all powers,
friendly and unfriendly, understand by what is
right only what self-interest tells them about their own
interest and honor?'' The politician knows that ``the world is not
governed with the Lord's Prayer.'' The bitterness with which he
rebuffs the moralist is mixed with melancholy. Naturally
all would be only too glad to be out of the painful entanglements;
but there are meshes in the net of fate that
cannot be torn apart; and the net encloses all,
even the mightiest struggle in the toils.

Is there, then, no way out? Yes, certainly. But to
find it and choose it and follow it there is required
a self-reflection that can be attained only through
painful experiences.

The first question that presents itself for answer
is this: how did it actually come about that politics
and conscience have moved so far away from each other?
It is a well-known saying that he who wills the end must also
will the means. From this it has been inferred that the end
sanctifies the means, and from this in turn it has been inferred that bad means
may be used for the attainment of a good end. This
last is obviously Jesuitical. Yet it must be remarked
that the Jesuits were not the first inventors; the discrepancy is
of older date and of deeper origin.

% --- p. 91 ---
\setcounter{footnote}{0}%
\opage{91}The discrepancy between means and end can first be
recognized in its depth when the end is eternal and the means
temporal, the end holy but the means worldly.
The politics of the ancient world, especially of the Romans, was in
many respects immoral, arbitrary, revolting, and
cruel; and yet the unconscionableness in all this
coarseness and selfish unruliness is only child's play in
comparison with the ecclesiastical politics of the Middle Ages.
Before conscience can properly be violated, it must
have awakened, its ideal demand have been acknowledged, and a
universally valid standard have been established for the judgment of
actual deeds. But the conscience of the ancient world
was wrapped in natural self-conceit and
stirred only in mythical dreams. It is not merely
a Nero who is unconscionable; but even if ethics
turns to a Cato the Elder, rips him open, and
looks him over inside along the seams, it will hardly
find a trace of what we must call conscience. The
only people of antiquity of whom it could be said in the stricter sense
that they had conscience were the Jews; and yet the Jewish people is
known to have applied a good part of its
natural understanding to evading the Law and cheating
Jehovah. Only with Christianity does conscience awaken
in history.

The characteristic feature of Catholic church politics
is now easy to illuminate. The Church naturally had to impress upon
men the demands of conscience; the aim it
set itself was supramundane, its end holy.
Nevertheless this holy end was, according to the Catholic point of view,
to be furthered and attained by worldly means. Means and end are thus as
dissimilar as is at all possible. The more the Church
% --- p. 92 ---
\setcounter{footnote}{0}%
\opage{92}became involved in worldly affairs, the more its
politics developed. Christ was no politician; the
only worldly means he employed were suffering and
self-sacrifice; and the means corresponded to the end, for his
kingdom was not of this world. But when the Church had
set itself the lofty goal of making the supramundane worldly,
it could not content itself with such meager means.
For making the supramundane worldly it had in
earnest to think of using worldly means; it
had to engage in politics. When conscience then shrank
from the means but yet approved the end, ``the
holy end,'' it lay near at hand to address it in a practical
way and to remind it that everyone who wills the end necessarily
must will the means. If conscience found itself revolted
by the means that were employed, how comfortingly soothing
was then not the assurance that the Church had means of atonement,
that the end sanctifies the means, that one can with a good
conscience employ bad means for good ends.
Papal politics was supposed, according to its ideal, to be a
politics of conscience, its boldest means means of conscience;
but in reality it was so self-seeking,
so arbitrary and intriguing that ordinary princely politics
had difficulty enough in matching it in
unconscionableness: it was a glaring contradiction.

The next step was that conscience went out and
politics was left alone. This turning point was
stated with clear consciousness by the Florentine Nicòlo
Machiavelli (1469--1527). By taking part in
public life he acquired an early taste for politics and practiced,
as it seems, even in private relations the principle
that one must not be too particular about the means in attaining an end.
He led an unprejudiced,
% --- p. 93 ---
\setcounter{footnote}{0}%
\opage{93}unrestrained existence, was in love not only in his youth
but even in an advanced age, sang serenades,
wrote indecent comedies, and was to such a degree the good
head of his comrades that they, during his absence as
an exile, wrote to him: ``Now that you are not here
any more, no one understands cards, drink, or
any other merriment.'' From the Church he learned that he who
by cunning and power wishes to make his fortune on the great theater
must be a good actor. To try his powers
in devout roles he joined one of the
pious brotherhoods and, when his turn came,
delivered a penitential sermon that ended with a heart-crushing exhortation,
full of unction: ``We ought to imitate St. Francis
and St. Jerome, who, in order to mortify the flesh and rob it of
power to tempt them, rolled in thorns or struck
their breasts with a stone; we, on the contrary, are deceived by
pleasure, lulled in errors, entangled in the snares
of sin, and are in the hands of the Devil. Therefore we must,
in order to escape from it, do penance and cry with David:
miserere mei Deus, or weep with St. Peter.''

If religion was a church fraud, a stage play, what
then was morality? In his book, as famous as it is notorious,
\emph{The Prince} (il principe), Machiavelli gives an account of
political justice. ``The ruler shall always have on
his lips words like justice, lawfulness, mildness, and
religion, but not trouble himself about these things when
he has an advantage in the opposite. Cruelty is necessary
for a new government, and one shall rather make
oneself feared than loved, when one cannot become both
at once. The aim of a government is its duration,
and that cannot be attained without hardness, because
men are usually ungrateful, false, and
% --- p. 94 ---
\setcounter{footnote}{0}%
\opage{94}quarrelsome, so that they must be held in check by fear
of punishment. The Prince must therefore know how to use
both poison and dagger, both cunning and force; he must be able to
play both the fox and the lion.'' How seriously Machiavelli
means it, and how far his conception of
justice deviates from the Socratic one in Plato's Republic, is seen
among other things from ``The Discourses on the Decades of Livy,''
in which he often refers to \emph{The Prince}. He declares
openly that justice arose through men's
having seen how useful the good and how harmful
the evil was, that men let themselves be moved
to the good only by necessity and self-interest,'' and so on.
Remarkable indeed, that the principles which the ``heathen''
Plato puts in the mouth of an insolent
sophist are, precisely in a Catholic-Christian age, proclaimed
openly and to all the world by the ``penitential preacher'' Machiavelli!
It is a great transformation of consciousness, a yield of world-historical
significance, that we have before us here.

Machiavellianism has this in common with the political
view of the ancient world, that it sets the state as
the highest, the only thing that has validity in itself.
Greek and Roman history has Machiavellian
traits in word as well as in deed. That Machiavellianism
nevertheless differs from the heathen
view of the state follows from the fact that it has Christianity,
especially Catholicism, as its presupposition. When a heathen
sophist asserts that nothing is in and for itself good
or evil, he can yet in a way be excused on the ground
that the conscience he knows is a blurred
essence; but when a Machiavelli openly sets forth
the same assertion, then it happens with scorn % print: Paaatand (misprint)
and denial of the clearly recognized conscience.
% --- p. 95 ---
\setcounter{footnote}{0}%
\opage{95}When a statesman in antiquity regarded the individuals of society
as without rights over against the state and therefore did not
hesitate to sacrifice them by the thousands for some
political purpose, the revolting character of this is
softened by the notion that the personal in
human worth and human right was still veiled in mythical
fog; but when Machiavellianism, with unveiled
praise of falsehood, dissimulation, perjury, and fratricide,
coolly and with strict consistency develops the same
fundamental thought, then that means that the ideal of man's
personal worth has appeared in the light of day only
to be trampled underfoot with shameless insolence.

One may object that neither ecclesiastical nor
secular politics has ever avowed
Machiavellianism, still less carried through its consequences.
But what does such an objection mean? To
avow Machiavellianism honestly and openly would
surely be the stupidest thing a Machiavellian could do.
As for carrying it through, a Louis XI in
France and a Ferdinand the Catholic in Spain
may well be regarded as talented precursors in practice,
and as for the theory itself, it is a fact
that Catherine de' Medici, when she was married
to Henry II, brought the book on ``The Prince''
with her from Italy. That she studied it
thoroughly, and that the fruits of her studies have shown themselves in
a brilliant fashion, especially in
the wedding of blood at Paris (the night of St. Bartholomew 1572), is
well known. As long as secular politics had a
religious cast, and the hydrogen flame of fanaticism intensified
the heat, the temperature was in a high degree favorable to
the flowering of Machiavellianism.

% --- p. 96 ---
\setcounter{footnote}{0}%
\opage{96}The ecclesiastical made good the secular, when the prince
committed a tyrannical outrage, and the Pope's holy
hand washed off the blood stain (the Stockholm Bloodbath).

To what degree the Catholics had a sense for
Machiavelli's doctrine becomes clear to us when we see how
shrewdly, piously, and cunningly it was conceived and
further developed by the Jesuits: Jesuitism is, politically
understood, an ecclesiastically applied Machiavellianism. That
religious Machiavellianism, even without regard to
the ecclesiastical, has been appropriated by the Protestants proves
how deep roots it has in human
nature. When Charles X (Gustavus) had broken the Peace of Roskilde,
and with the intention of seizing Denmark-Norway
had landed with his fleet at Korsør, he struck up
the hymn: ``In Thee I hope, O Lord most dear!'' If
it can succeed for a prince to warp his Protestant conscience
to such a degree that God in heaven himself becomes
a Machiavellian, then he can also, with Charles Gustavus,
without blushing speak of his conquests as ``gifts he
owes to God's grace.'' In the spirit of this politics Charles
Gustavus's confidential friend Schlippenbach has in confidence
communicated to Elector Frederick William of Brandenburg an
edifying interpretation of the revelation of grace in politics:
``God no longer reveals his will to princes
through prophets or dreams; but when there
shows itself a good opportunity to harm one's neighbor, it ought
to be held a divine calling.'' --- The Great
Elector without doubt felt so edified by the
interpretation that he saw no obstacle
to a sincere union between Lutherans and Reformed.

% --- p. 97 ---
\setcounter{footnote}{0}%
\opage{97}True, in literature there gradually appeared a
number of moral refutations of Machiavelli's doctrine; but
practice was stronger than theory. That holds not
least of the Brandenburg practice. Frederick II of
Prussia wrote in his youth an Anti-Machiavelli, a
work full of virtuous indignation over the immoral principles;
but the King of Prussia, whose political conscience
on the whole seems to have been in a
somewhat mythical state, came precisely by this youthful work
to break the staff over principles by whose
clever application he was afterwards to shine in history as the conqueror of Silesia
and the hero of the Seven Years' War. This
inconsistency is, from a Machiavellian point of view,
for the rest quite consistent.

In our century Machiavellianism may on the whole
be said to be on the decline. Only now and
then has the immoral in politics come to a breakthrough
in a revolting manner. In order not to come too near
the present and tear open the wounds that still bleed,
we shall limit ourselves to recalling the English
attack on Copenhagen in 1807. This step was for
natural feeling so revolting that a Danish poet
has even wished to make clear to us what disapproval
it must have aroused in the English people itself.
Well known is the song of the English captain: ``with
zest I cried: God and King George! Every time I
fired at Napoleon's brigs, But here I almost
must weep with sorrow: On robbery now our fleet lies.
--- Sail on, you! I obey not the King's command, But go
on other adventures and voyages, It stands in Scripture:
You shall obey God More even than you shall obey Rome's
Emperor!''

% --- p. 98 ---
\setcounter{footnote}{0}%
\opage{98}Let it be with the legend as it will; certain it is that
English politics did not show itself unworthy of Machiavellianism.
For so great an errand men of the right school
were required. The English ministry had in
Sir Brook Taylor and Sir Francis Jackson precisely two such
men of honor. Sir Brook Taylor already had
a name in the great world. ``He had namely formerly,
when he represented England at Kassel, taken part
in the attempts to murder First Consul Bonaparte'';
the end sanctifies the means. As to Sir Francis
Jackson's prominent personality, one of
our younger historians\footnote{Alex. Thorsøe Den danske Stats politiske Historie. Første Tidsrum. Copenhagen 1873. P. 154 ff.} has furnished a characterization that
deserves to be read. As a \emph{Deus} machina (a god out of % print: Deus machina (sic)
a machine) this diplomatic agent appears suddenly
at Kiel and reveals to Crown Prince Frederick England's
hidden plans. He ``began his message with
the declaration that the English government had procured
reliable evidence that Napoleon
would force Denmark to close her ports to the English
and to take part in the Continental Blockade.
Both Great Britain's interest and security and
Denmark's independence demanded that the court at Copenhagen
enter into an intimate relation with England and consent
that its fleet be brought to England before
the frost had shut it in in the harbor, so that it
might be protected against Napoleon's pursuit.''
He solemnly assured, ``that the fleet should be
given back when peace was concluded.'' He offered
the Crown Prince a close alliance with Great Britain,
% --- p. 99 ---
\setcounter{footnote}{0}%
\opage{99}guarantees, and so on, and added finally, ``that if the
Danish government by such a step feared to
bring itself into a dangerous position toward France, it
must yet not forget that the English force
which now lay before Copenhagen was so considerable that it
might well look as if it had yielded to
compulsion, to which England for the rest would resort, if Denmark
did not enter into its plans.'' --- Proposals of so
pithy a content had naturally to be put forward in a
correspondingly pithy form, and Sir Francis Jackson was a
master of the pithy form. It is reported, ``that Jackson
in the beginning said nothing but ``De'' to the Crown Prince,
whereupon the latter asked whether he did not know
whom he was speaking with, that it was a prince,
the son of the King of Denmark, whereupon Jackson became somewhat
more polite'' and so on.

The extreme brutality with which the apostles of Machiavellianism
presented themselves on this occasion
had, unfortunately, the result that the worthy representative of honesty
and morality, the noble Crown Prince, was
not able to come to political reflection before
it was too late. ``Crown Prince Frederick ought certainly
to have sought to gain time by entering into negotiations
with Jackson, until the intentions of the French government
had become known to him. He would then, when this
had happened, have come to the recognition that the Emperor's
plans with regard to Denmark were just as
contrary to international law and just as little reconcilable
with the free right of self-determination of the Danish government
as those of the British cabinet. The Crown Prince would then have
come to know by experience that Napoleon, with the concentration of troops
at the frontier of Holstein, quite connected the same
% --- p. 100 ---
\setcounter{footnote}{0}%
\opage{100}thought as the English ministry with the gathering
of the fleet in the Sound, and it would thus have stood
perfectly clear to him that the two powers were in equal degree
guided by self-interest, and that both with equal
ruthlessness and disregard of Denmark's freedom of action
sought only through force and compulsion to carry through their
own will. But Crown Prince Frederick did not
enter into negotiations with Jackson at all, and if
one cannot admire his ``diplomatic acumen,'' yet
one receives a so much the more satisfying impression of his
moral character. ``He declared---which in
that time, when political acts of violence belonged to the order
of the day, sounded strange---that there was hardly in history
any example of an attack like that with which
they dared to threaten Denmark, he stated that among
the barbarian states one might expect to find greater honesty
than with the English government, on whose help
one could not even rely, even if it were an ally.
When Sir Francis Jackson lodged an objection to
this statement, Crown Prince Frederick remarked that the
man who had the moral courage to undertake a
mission like his must also be able to bear to hear
everything; when the English envoy assured that England
would compensate Denmark with money for the damage suffered,
he was interrupted with the famous words: ``And
with what will one compensate Denmark's honor?''

Machiavellianism is not limited to certain
forms of government; it thrives about equally well in them
all. --- The democrats would gladly have us believe that when
the majority comes to the helm, the last hour of the great injustices
will strike, and Politics and Conscience
fraternally embrace each other; for under
% --- p. 101 ---
\setcounter{footnote}{0}%
\opage{101}majority rule means and ends are to be dealt with honestly.
How beautiful that is, if one could only believe it!
But alas, history teaches that Machiavellianism
has precisely in the power of the majority an old ally,
and that is party spirit. ``He who wills the end must
will the means; the means is good when it leads to the
goal.'' The party makes individuals into means, uses
them as long as they serve its end, and throws them
away when it no longer needs them. Conversely
the participants in turn make the party a means, and leave
it when they have used it as a footstool. It is
just the liberal, the genuinely free-minded thing about party spirit
that its claims upon the mutual agreement of the participants
are by no means importunate. It is far from
demanding that the members of the party be agreed
on anything so essentially inward that it could trouble
conscience; it is enough that they are agreed on some
external matter, such as overthrowing a ministry, helping
their leaders to power, taking revenge on opponents,
and the like. Every such external end is as
a rule only a means to other and higher ends.
Everyone in the party can then freely go about and fawn
on his conscience and persuade it that the petty tricks
which the party now and then permits itself are only innocent knacks
by which one paves a way for the spirit. True, there are
many in the party who as yet do not have the spirit, but
the party is, after all, only a transition, a means. True,
some of the means the party employs are perhaps not
quite lawful, and would hardly deserve the approval of
conscience if it were the opponents who used them;
but one may---in view of the political situation---
allow oneself a good deal, when it happens for a good end.
% --- p. 102 ---
\setcounter{footnote}{0}%
\opage{102}``The opponents have---fortunately---no good end.
There is in what they say and do, to be sure, something that we
would find right if it were we who said it,
we ourselves who did it, but from the side of the opponents it is
---on account of the political situation---nonetheless
unreasonable and dishonorable; for the fact is that the party's
opponents without exception are idiots.'' By
such reflections a decent man can easily get his
conscience chloroformed, even to such a degree that even
the boldest operations with the political knife do not
cause it pain. This is the most questionable kind.
of Machiavellianism.

When Hercules, the myth relates, on his western expedition
came to Libya, he encountered the giant Antaeus, whose
cruelty to strangers made him a veritable
monster. Antaeus forced the newcomers to wrestle
with him, and when he had overcome them, he strangled
them and set up their skulls around his
dwelling. What made him invincible in the struggle
was his mother, \emph{Earth}, whereby his strength was doubled.
But Hercules seized him with his arms around the
waist and held him suspended in the \emph{air}, until he,
strangled by the hero's strength, gave up his spirit. --- Machiavellianism
is like the giant Antaeus; it cannot be defeated
by merely worldly forces; for worldliness is
its mother, its innermost essence is the selfishly worldly.
If the monster is to be smothered, then it must be lifted up into the sphere of
religion and morality, into the atmosphere
in which conscience breathes, and held up there by mighty
arms. No diplomatic artist and no political
genius is able to hold out against the monster; it
blinds the clever with their own cleverness and smothers
% --- p. 103 ---
\setcounter{footnote}{0}%
\opage{103}the strong in foul vapors. If there is to be any
power on earth that is able to seize the monster about the
waist and hold it suspended in the air, then it must be
the popular spirit, when it truly awakens. The popular spirit is
a Hercules of moral forces, a real spirit in a
spiritually active people. It is a hero who goes
honestly to work; on his banner stands humanity's
demand in flaming script, and the cunning of worldliness
and the deviousness of politics must, ashamed, bow
beneath it. If he comes on his expedition to
the wrestling match, then he will surely use honest, if
not exactly lenient means: he will subdue the petty
defiance, air out the camps of the parties and cleanse the encumbered
state, as his mythical predecessor of old
cleansed the stable of Augeas.

\end{document}
